\documentclass[11pt,oneside,openany]{book}
\usepackage[T1]{fontenc}
\usepackage{lmodern}
\usepackage[letterpaper,margin=1.10in,includeheadfoot]{geometry}
\usepackage{amsmath,amssymb,amsthm,mathtools}
\usepackage{microtype,xcolor,titlesec,fancyhdr,etoolbox,needspace}
\usepackage{hyperref}

\ifdefined\pdftrailerid
  \pdftrailerid{MathGenome Topics in Field and Galois Theory}
\fi
\definecolor{ink}{HTML}{172923}
\definecolor{deepgreen}{HTML}{2B5242}
\definecolor{rulegreen}{HTML}{8A9F92}
\hypersetup{colorlinks=true,linkcolor=deepgreen,urlcolor=deepgreen,citecolor=deepgreen,
  pdftitle={Topics in Field and Galois Theory},pdfauthor={MathGenome},
  pdfsubject={ALG-04-D01 through ALG-04-D04},bookmarksnumbered=true}
\setcounter{tocdepth}{1}
\setlength{\headheight}{15pt}
\setlength{\emergencystretch}{2em}
\newcommand{\runningchapter}{}
\raggedbottom
\titleformat{\chapter}[display]
  {\normalfont\color{ink}}{\color{deepgreen}\large\scshape Chapter \thechapter}{8pt}
  {\huge\bfseries}
\titlespacing*{\chapter}{0pt}{7pt}{22pt}
\titleformat{\section}{\normalfont\Large\bfseries\color{ink}}{\thesection}{0.7em}{}
\pagestyle{fancy}
\fancyhf{}
\fancyhead[L]{\small\sffamily\color{deepgreen}MathGenome}
\fancyhead[R]{\small\nouppercase{\runningchapter}}
\fancyfoot[C]{\small\thepage}
\renewcommand{\headrulewidth}{0.35pt}
\renewcommand{\chaptermark}[1]{\renewcommand{\runningchapter}{\chaptername\ \thechapter}}
\fancypagestyle{plain}{\fancyhf{}\fancyfoot[C]{\small\thepage}
  \renewcommand{\headrulewidth}{0pt}}

\newtheorem{theorem}{Theorem}[chapter]
\newtheorem{lemma}[theorem]{Lemma}
\newtheorem{proposition}[theorem]{Proposition}
\newtheorem{corollary}[theorem]{Corollary}
\theoremstyle{definition}
\newtheorem{definition}[theorem]{Definition}
\newtheorem{example}[theorem]{Example}
\theoremstyle{remark}
\newtheorem{remark}[theorem]{Remark}
\newcommand*\keeptheoremheading[1]{\BeforeBeginEnvironment{#1}{\Needspace{4\baselineskip}}}
\forcsvlist{\keeptheoremheading}{theorem,lemma,proposition,corollary,definition,example,remark}
\allowdisplaybreaks[1]

\begin{document}
\frontmatter
\hypersetup{pageanchor=false}
\begin{titlepage}
\thispagestyle{empty}
\thispdfpagelabel{Cover}
\begin{flushleft}
\vspace*{0.72in}
{\large\scshape\color{deepgreen}MathGenome\par}
\vspace{0.36in}
{\color{rulegreen}\rule{0.94\textwidth}{0.65pt}\par}
\vspace{0.46in}
{\color{ink}\fontsize{35}{41}\selectfont\bfseries
  Topics in Field and\\[2pt]Galois Theory\par}
\vspace{0.34in}
{\color{deepgreen}\Large Picard--Vessiot theory, Galois cohomology,\\[4pt]
  and embedding problems\par}
\vfill
{\normalsize ALG-04-D01 through ALG-04-D04\par}
\vspace*{0.50in}
\end{flushleft}
\end{titlepage}
\hypersetup{pageanchor=true}
\begingroup
\small
\makeatletter
\patchcmd{\l@chapter}{\vskip 1.0em \@plus \p@}
  {\vskip 0.45em \@plus 0.25\p@}{}{}
\makeatother
\tableofcontents
\endgroup

\chapter*{Scope and conventions}
\phantomsection
\addcontentsline{toc}{chapter}{Scope and conventions}
\renewcommand{\runningchapter}{Scope and conventions}

These notes continue the field theory and profinite group theory developed
in \emph{Infinite Galois Theory} and \emph{Topics in Group Theory}.
The opening chapters study linear systems over differential and difference
fields. We construct their Picard--Vessiot rings and explain how
automorphisms of a solution ring determine intermediate objects. The
total quotient ring matters in the difference setting: even over a field,
a solution ring can have several connected components.

We then study the continuous cohomology of the absolute Galois group with
discrete coefficients. The comparison with finite quotients gives concrete
descriptions of classes by field elements and by finite extensions.
The last chapter relates extensions of finite groups to lifting problems
for the absolute Galois group and proves explicit specialization results.

Chapter 4 proves specialization at a specified separable fibre and
full-group specialization in explicit families. For general
embedding problems, see Malle--Matzat~\cite[Chapter IV, Sections 1--2]{mallematzat}.
For Hilbert irreducibility, regular realization and rigidity, see
D\`ebes~\cite[Sections 5.2--5.3 and 8.2]{debes}.
Haran--Jarden~\cite[Sections 1 and 5]{haranjarden} treats regular split
embedding problems over complete valued fields.

Unless another convention is specified, actions are on the left and
\(st\) acts by first applying \(t\), then \(s\). A fixed separable
closure \(k^{\mathrm s}\) of a field \(k\) carries the discrete topology
when used as a coefficient group, and
\(\operatorname{Gal}(k^{\mathrm s}/k)\) carries the Krull topology.
The first two chapters state their characteristic and constant-field
assumptions separately from those used in ordinary Galois theory.

Proofs give the hypotheses for each use of earlier MathGenome notes.
The bibliography identifies the source literature, and the accompanying
map locates each numbered statement in the source and the PDF and records
its dependencies.

\mainmatter
\chapter{Picard--Vessiot Rings}

For an algebraic polynomial one adjoins roots. For a linear differential
or difference system one adjoins an invertible matrix of solutions,
and the role of the splitting field is played initially by a ring.
This chapter explains why the difference ring can have zero divisors,
why no new constants appear under the stated hypotheses, and how to
recognize the resulting extension. The two theories will use the same
formalism in the next chapter.

\section{The operator and the system}

\begin{definition}[Operator rings and constants]
\label{res:PV01}
All rings in these chapters are commutative and have an identity. In
the differential case an operator is a derivation \(\partial\), an
additive map satisfying \(\partial(ab)=a\partial b+b\partial a\).
In the difference case it is a ring automorphism \(\sigma\).
An operator subring is closed under \(\partial\), or under both
\(\sigma\) and \(\sigma^{-1}\), respectively. An operator ideal is
defined by the same stability condition. An operator ring is
\emph{simple} if it is nonzero and its only operator ideals are \(0\)
and the whole ring. Its constants are the kernel of \(\partial\), or
the fixed ring of \(\sigma\), respectively, and will be denoted
\(S^{\mathrm{const}}\) for either kind of operator on \(S\).
An operator field is a field with the indicated operator.
\end{definition}

Except where explicitly stated, \(k\) is an operator field of
characteristic zero and \(C=k^{\mathrm{const}}\) is algebraically
closed. On a tensor product of extensions over \(k\) the differential
operator is the sum of the operators in the two factors and the
difference operator acts in both factors. The product rule and the
common operator on \(k\) make these formulas well defined.
Finite-type algebras over a noetherian base are noetherian
by the Hilbert basis theorem in~\cite[Section 23.11, PDF p.~191]{mg-d07}.

\begin{lemma}[Fundamental matrices and scalar equations]
\label{res:PV02}
In the differential case let \(A\in M_n(k)\) and consider
\(\partial Y=AY\). In the difference case let \(A\in GL_n(k)\)
and consider \(\sigma(Y)=AY\). An invertible solution matrix is
called fundamental. If \(Y\) is fundamental in an extension \(S\),
every solution column in \(S^n\) is uniquely of the form \(Yc\) with
\(c\in(S^{\mathrm{const}})^n\). Two fundamental matrices in \(S\)
differ on the right by an invertible constant matrix.
A monic scalar equation of order \(n\) gives such a system by taking
the column consisting of the unknown and its first \(n-1\) derivatives
or iterates. For the difference equation its order-zero coefficient
must be nonzero to give an invertible system matrix.
\end{lemma}

\begin{proof}
For a solution column \(v\), differentiating \(Y^{-1}v\) gives
\(-Y^{-1}Av+Y^{-1}Av=0\). In the difference case
\(\sigma(Y^{-1}v)=(AY)^{-1}Av=Y^{-1}v\).
The converse follows by applying the operator to \(Yc\).
Apply this to the columns of a second fundamental matrix and then
interchange the two matrices to obtain constant inverses.
For the scalar equation, the first \(n-1\) rows of its system shift the
entries of the column by one. The last row is the negative of the
coefficients in increasing order. In the difference case its determinant
is the order-zero coefficient times a sign, so the stated condition
is precisely invertibility.
\end{proof}

\section{Simplicity, constants and the total ring}

\begin{lemma}[The components of a simple operator ring]
\label{res:PV03}
A simple noetherian differential ring containing \(\mathbb Q\) is a
domain. A simple noetherian difference ring is a finite product of
domains whose factors are permuted transitively by the automorphism.
In either case the operator extends to the total ring of fractions,
obtained by inverting the non-zero-divisors. This ring is a field in
the differential case and a finite product of fields with a transitive
permutation in the difference case. A simple ring and its total ring
of fractions have the same constants, and these constants form a field.
\end{lemma}

\begin{proof}
If \(x^m=0\), applying a derivation \(m\) times shows that
\(m!(\partial x)^m\in(x)\), because the terms other than the one
with one derivative in each factor have an undifferentiated factor.
Thus the nilradical is a differential ideal in characteristic zero
when the ring contains \(\mathbb Q\). For every minimal prime \(P\),
localize at \(P\) and apply this argument to the nilradical \(PR_P\).
Contraction to the original ring shows that \(P\) is a differential
ideal. Simplicity gives \(P=0\), so the ring is a domain.

In the difference case the nilradical is a proper stable ideal, so
the ring is reduced. The automorphism permutes its finite set of
minimal primes. To recall finiteness, a minimal counterexample
among closed sets of a noetherian spectrum would split into two
smaller closed sets, each a finite union of irreducible sets.
Applied to the whole spectrum, this gives finitely many
irreducible components and hence finitely many minimal primes.
The intersection of any proper subset of the
minimal primes is nonzero: choose an omitted minimal prime \(P\)
and multiply elements, one in each prime of the subset and outside
\(P\). This also applies to the empty subset, using the product \(1\).
Simplicity therefore implies transitivity. For the minimal primes
\(P_i\), the sum of the ideals \(J_i=\bigcap_{j\ne i}P_j\) is nonzero
and is stable, so it equals the ring. Write \(1=\sum_i e_i\) with
\(e_i\in J_i\). The image of \(e_i\) modulo \(P_j\) is \(1\) for
\(i=j\) and \(0\) otherwise. The intersection of the \(P_j\) is zero,
so the \(e_i\) are orthogonal idempotents. The components \(e_iS\)
are the domains \(S/P_i\). A non-zero-divisor has a nonzero image
in each component, which proves the assertion about the total
ring of fractions.

The quotient rule extends \(\partial\), and
\(\sigma(r/s)=\sigma(r)/\sigma(s)\) extends the automorphism.
A nonzero constant of a simple ring generates a nonzero stable
ideal, hence is a unit with constant inverse. If \(q\) in its total
ring of fractions is constant, the ideal
\(\{a\in S:aq\in S\}\) is stable and contains a denominator of \(q\).
It therefore contains \(1\), so \(q\) was already a constant of \(S\).
\end{proof}

\begin{lemma}[Specialization in a finitely generated domain]
\label{res:PV04}
Suppose that \(B\) is a finitely generated domain over an infinite
field \(F\) and \(t\in B\) is transcendental over \(F\).
Then \(t-a\) is a nonunit of \(B\) for all but finitely many \(a\in F\).
Consequently a field finitely generated as an algebra over \(F\)
is finite algebraic over \(F\). If \(F\) is also algebraically
closed, a nonzero reduced finitely generated \(F\)-algebra
is detected by its \(F\)-algebra maps to \(F\): every nonzero
element is nonzero under at least one of them.
\end{lemma}

\begin{proof}
Extend \(t\) to a transcendence basis \(t,u_1,\dots,u_j\) of
\(\operatorname{Frac}(B)/F\) by choosing some of the generators of
\(B\). Each generator satisfies a monic polynomial over
\(F(t,u_1,\dots,u_j)\). Choose a single nonzero polynomial \(d(t,u)\)
clearing all the denominators in their coefficients. Then
\(B[d^{-1}]\) is integral over
\(D=F[t,u_1,\dots,u_j,d^{-1}]\).
For all but finitely many \(a\in F\), \(d(a,u)\ne0\):
view the coefficients in the \(u_i\) as polynomials in \(t\) and take
one nonzero coefficient. At such \(a\), specialization at \(t=a\)
exhibits \(t-a\) as a nonunit in \(D\). It is still a nonunit in
\(B[d^{-1}]\): an inverse integral over \(D\) would belong to \(D\),
by multiplying its monic relation by a suitable power of \(t-a\).
It follows that \(t-a\) is a nonunit in \(B\).

A finitely generated \(F\)-algebra that is a field can therefore
have no transcendental element, and finitely many algebraic
generators give a finite extension. Over algebraically closed \(F\),
a maximal ideal in a nonzero finitely generated \(F\)-algebra has
residue field \(F\). To detect a nonzero element \(b\) in a reduced
algebra, apply this to the nonzero finitely generated algebra
obtained by inverting \(b\), and restrict the resulting map.
\end{proof}

\begin{proposition}[No new constants in a simple finite-type algebra]
\label{res:PV05}
Let \(S\) be a simple operator ring finitely generated as an algebra
over \(k\). Then \(S^{\mathrm{const}}=C\); the same holds for
its total ring of fractions.
\end{proposition}

\begin{proof}
By Lemma~\ref{res:PV03}, the simple ring is a finite product of
domains, with one factor in the differential case. Suppose a
constant \(c\in S\) is transcendental over \(k\). Some component
of \(c\) is transcendental, since otherwise the product of
polynomial relations in the finitely many components would give a
relation in \(S\). For every \(a\in C\) the difference \(c-a\) is a
nonzero constant, hence a unit in \(S\). Its component in the
chosen domain is a unit as well. As \(C\) is infinite, this
contradicts Lemma~\ref{res:PV04} with \(F=k\).

Thus \(c\) has a monic relation of least degree in \(k[T]\).
In the differential case, applying \(\partial\) to that relation
gives a relation of lower degree with coefficients the derivatives of
the original coefficients. In the difference case, applying
\(\sigma\) and subtracting the original relation gives one whose
coefficients are the differences of the coefficients and their
images. In both cases the lower-degree relation is zero, hence the
original coefficients lie in \(C\). Since \(S^{\mathrm{const}}\)
is a field and \(C\) is algebraically closed, \(c\) belongs to
\(C\). The statement about the total ring follows from
Lemma~\ref{res:PV03}.
\end{proof}

\section{Construction and uniqueness}

\begin{definition}[Picard--Vessiot ring and its total extension]
\label{res:PV06}
For either system of Lemma~\ref{res:PV02} a Picard--Vessiot ring
over \(k\) is a simple operator extension \(R\) containing a
fundamental matrix \(Y\) and generated as a \(k\)-algebra by its
entries and \(\det(Y)^{-1}\). Its total extension \(Q(R)\) is the
total ring of fractions with the extended operator. It is called a
Picard--Vessiot field when it is a field.
\end{definition}

\begin{theorem}[Existence with prescribed constants]
\label{res:PV07}
Every system of Lemma~\ref{res:PV02} over \(k\) has a
Picard--Vessiot ring. Both the ring and its total extension have
exactly the constants \(C\).
\end{theorem}

\begin{proof}
Let \(U=k[X_{ij},\det(X)^{-1}]\) with \(X\) an indeterminate
square matrix. A derivation of the polynomial ring is defined by
\(\partial X=AX\) and extends to the localization by the
quotient rule. In the difference case set \(\sigma(X)=AX\) and
apply \(\sigma\) to the coefficients. The image of \(\det X\) is
\(\det A\det X\), a unit in \(U\). This is an automorphism
with inverse acting on coefficients by \(\sigma^{-1}\) and on \(X\)
by \(\sigma^{-1}(A^{-1})X\).
By the ascending chain condition, a maximal proper operator
ideal exists, and its quotient is simple. A proper ideal does
not meet \(k\) nontrivially, so the quotient is an extension
of \(k\). The image of \(X\) is fundamental and generates the
required algebra. Proposition~\ref{res:PV05} supplies its constants.
\end{proof}

\begin{lemma}[Descent of ideals with constant variables]
\label{res:PV08}
Let \(S\) be a simple operator ring with constant field \(C\).
If \(D\) is any \(C\)-algebra, equip \(S\otimes_C D\) with the
operator from the first factor only. Every operator ideal is extended
from its contraction to \(D\), and the constants equal \(1\otimes D\).
In particular, if an operator extension \(B\) of \(S\) is generated
as an \(S\)-algebra by a \(C\)-subalgebra \(D\subseteq
B^{\mathrm{const}}\), multiplication gives an isomorphism
\(S\otimes_C D\cong B\).
\end{lemma}

\begin{proof}
For an operator ideal \(I\), divide \(D\) by its contraction and
replace \(I\) by its image. Exactness of tensoring over \(C\) gives
zero contraction for this image. Suppose the image contains a
nonzero \(\sum_{i=1}^m s_i\otimes d_i\) with \(m\) minimal,
and choose the \(d_i\) linearly independent over \(C\).
The possible first coefficients in elements of \(I\) whose
expressions use only \(d_1,\dots,d_m\) form a nonzero operator
ideal of \(S\). It is the whole ring. Thus we can choose such
an element with first coefficient \(1\). Applying \(\partial\),
or subtracting the element from its image under \(\sigma\), now
gives an element of \(I\) with at most \(m-1\) terms. By minimality
it is zero, so all the coefficients are constants. This gives
a nonzero element of the contraction, a contradiction.
Expanding an arbitrary element of the tensor product in a finite
linearly independent set in \(D\) also shows directly that a
constant has constant coefficients. The multiplication map in
the last assertion is surjective with an operator kernel
whose contraction to \(D\) is zero. The first assertion gives
injectivity.
\end{proof}

\begin{theorem}[Uniqueness and recognition inside a simple extension]
\label{res:PV09}
Any two Picard--Vessiot rings for the same system over \(k\) are
isomorphic over \(k\) as operator rings, and the isomorphism
extends to their total extensions. More generally, if a simple
operator extension \(S/k\) has constants \(C\) and contains a
fundamental matrix \(Y\), the subalgebra
\(k[Y_{ij},\det(Y)^{-1}]\) is a Picard--Vessiot ring.
\end{theorem}

\begin{proof}
Take a Picard--Vessiot ring \(R'\) with fundamental matrix \(Y'\).
In \(S\otimes_k R'\) put
\(Z=(Y\otimes1)^{-1}(1\otimes Y')\). Its entries and inverse
determinant are constant by Lemma~\ref{res:PV02}.
The tensor product is nonzero, contains the first factor faithfully,
and is generated over that factor by
\(D=C[Z_{ij},\det(Z)^{-1}]\). Apply Lemma~\ref{res:PV08}
and choose a \(C\)-algebra map \(D\to C\) by
Lemma~\ref{res:PV04}. Evaluation defines an operator map
\(R'\to S\) taking \(Y'\) to \(YB\) for some \(B\in GL_n(C)\).
It is injective by simplicity of \(R'\), and its image is
exactly the algebra generated by \(Y\) and its inverse
determinant, since \(B,B^{-1}\) have entries in \(k\).
This proves recognition. Taking \(S\) itself to be a second
Picard--Vessiot ring gives uniqueness, and localization gives
the assertion about the total extensions.
\end{proof}

\begin{example}[A difference Picard--Vessiot ring that is not a field]
\label{res:PV10}
Let \(k=C(z)\), where \(C\) is algebraically closed of
characteristic zero, and set
\(\sigma(f)(z)=f(qz)\) for a fixed \(q\in C^\times\) of infinite
multiplicative order. The constants are \(C\). For
\(\sigma(y)=-y\), a Picard--Vessiot ring is
\(k[u]/(u^2-1)\) with \(\sigma(u)=-u\). It is already its
own total ring of fractions and has no embedding into a field.
\end{example}

\begin{proof}
A rational function fixed by the substitution \(z\mapsto qz\)
has its zeros and poles invariant under multiplication by \(q\).
Any orbit of a point in \(C^\times\) is infinite, so a nonzero
fixed rational function is \(az^m\) with \(a\in C^\times\) and
\(m\in\mathbb Z\). Invariance gives \(q^m=1\), so \(m=0\).
The idempotents \((1+u)/2,(1-u)/2\) exhibit the asserted ring
as \(k\times k\). The automorphism interchanges the two
idempotents, so no proper nonzero ideal is stable. The
invertible element \(u\) is fundamental and generates the ring,
which is therefore Picard--Vessiot. Its two orthogonal nonzero
idempotents preclude an embedding into a field.
\end{proof}

\paragraph{Sources.}
The differential construction follows the subject of
van der Put--Singer~\cite[Chapter 1, Sections 1.1--1.3]{vps-differential}
and Singer~\cite[Sections 1.2--1.3]{singer}. For the difference
construction see van der Put--Singer~\cite[Chapter 1, Section 1.1]{vps-difference}
and Roques~\cite[Definitions 7--8, Propositions 9 and 12, Theorems 13 and 17]{roques}.
The operator-ring viewpoint also connects to the treatment of Ore
operators in~\cite[Chapter 4]{mg-rm03d}.
\chapter{Linear Operator Galois Correspondence}

Fix a Picard--Vessiot ring \(R/k\), its fundamental matrix \(Y\),
and its total extension \(Q=Q(R)\), under the hypotheses of the
previous chapter. In this chapter an automorphism of an operator
ring always commutes with the operator. The difference case requires
intermediate total rings because even the simplest nontrivial
example in the preceding chapter can have orthogonal idempotents.

\section{The algebra of the Galois group}

\begin{lemma}[Reduced operator tensors in characteristic zero]
\label{res:PV11}
If \(A\) is a reduced finite-type algebra over a characteristic-zero
field \(F\), and \(L/F\) is any field extension, then
\(L\otimes_F A\) is reduced. The same conclusion holds when the base
is a finite product of such fields, \(A\) is again reduced and
of finite type, and \(L\) is a finite product of field extensions
of the base components.
\end{lemma}

\begin{proof}
The embeddings of the domains \(A/P\), with \(P\) running over the
finitely many minimal primes, into their fraction fields embed \(A\)
in a finite product of finitely generated fields over \(F\).
Tensoring over a field preserves this injection. For one such
fraction field \(M\), choose a transcendence basis \(t\) of
\(M/F\). Then \(M/F(t)\) is finite separable, so it is generated
by a root of a monic separable polynomial \(p\) by the primitive
element theorem~\cite[Section 14.31, PDF p.~113]{mg-d07}. The algebra
\(L\otimes_F F(t)\) embeds in \(L(t)\), as it is the localization
of \(L[t]\) at the nonzero polynomials of \(F[t]\).
The algebra \(L\otimes_F M\) is obtained from that localization
by adjoining a root of \(p\). The basis of powers of the root
exhibits an injection into \(L(t)[T]/(p)\), which is reduced
as \(p\) remains separable. This proves the field case. Decomposing
the tensor product componentwise proves the product case.
\end{proof}

\begin{theorem}[The torsor of a fundamental matrix]
\label{res:PV12}
Write \(Y_1=Y\otimes1\), \(Y_2=1\otimes Y\) in \(R\otimes_k R\),
and let \(Z=Y_1^{-1}Y_2\). Then \(Z\) is constant and the
constant algebra \(D=C[Z_{ij},\det(Z)^{-1}]\) is reduced. There
is a natural isomorphism
\[
 R\otimes_C D \ \xrightarrow{\ \sim\ }\ R\otimes_k R
\]
by using \(R\) in the first factor. The evaluations \(D\to C\)
identify the matrices of \(Z\) with the algebraic subgroup
\[
 G=\{g\in GL_n(C):Y\mapsto Yg
                   \text{ gives an operator automorphism of }R/k\}.
\]
In particular \(D=C[G]\). Every such automorphism extends uniquely
to \(Q\), and every operator automorphism of \(Q/k\) restricts to
\(R\). Replacing \(Y\) by another fundamental matrix conjugates
the matrix realization of \(G\).
\end{theorem}

\begin{proof}
Constancy of \(Z\) follows from Lemma~\ref{res:PV02}. Since
\(Y_2=Y_1Z\), the right tensor product is generated by \(D\)
over its first factor, and Lemma~\ref{res:PV08} proves the
isomorphism. The embedding of \(R\) into the product of its
fraction fields and Lemma~\ref{res:PV11} show that
\(R\otimes_k R\) is reduced, hence so is \(D\).

An evaluation \(D\to C\) sends \(Z\) to an invertible matrix \(g\);
it maps the second copy of \(R\) to the first by \(Y\mapsto Yg\).
This is an operator homomorphism, injective by simplicity and
surjective since \(Y=(Yg)g^{-1}\). Conversely any automorphism
determines an evaluation by sending \(a\otimes b\) to \(a\,g(b)\),
where \(g(b)\) denotes its action. Its matrix is constant by
Lemma~\ref{res:PV02}. Thus \(G\) is the zero set in \(GL_n(C)\)
of the relations of \(Z\), it is a subgroup by composition,
and its coordinate algebra is exactly \(D\) by
Lemma~\ref{res:PV04}. The result is the torsor formula: in
\(R\otimes_k R\) the second matrix is obtained from the first by
right multiplication by \(Z\).

An automorphism of \(R\) preserves non-zero-divisors and
extends uniquely. For an automorphism of \(Q/k\), the ratio of its
image of \(Y\) and \(Y\) is constant. Since \(Q^{\mathrm{const}}=C\)
by Theorem~\ref{res:PV07}, the image of \(Y\) is of the form \(Yg\)
and generates \(R\) over \(k\). This proves restriction. For
another fundamental matrix \(Yb\), the new coordinate matrix of
an automorphism is \(b^{-1}gb\).
\end{proof}

\section{Intermediate rings and invariant elements}

\begin{lemma}[Structure of an intermediate total ring]
\label{res:PV13}
Let \(F\subseteq Q\) be an operator subring containing \(k\)
such that every non-zero-divisor of \(F\) is a unit of \(F\).
In the differential case \(F\) is a field. In the difference
case \(F\) is a finite product of fields with the operator permuting
its factors transitively. In either case \(F^{\mathrm{const}}=C\).
If \(R_F=F[Y_{ij},\det(Y)^{-1}]\) inside \(Q\), then \(R_F\)
is noetherian and its total ring of fractions is \(Q\).
\end{lemma}

\begin{proof}
If \(Q\) is a field, the subring \(F\) is a domain and the
unit condition makes it a field. For the difference case write
\(Q=\prod_i Q_i\) by Lemma~\ref{res:PV03}. The distinct kernels
of the projection maps \(F\to Q_i\) are permuted transitively.
No two distinct such kernels are comparable: any strict inclusion
would, by a power of the operator carrying one kernel to the other,
lead to a cycle of strict inclusions in a finite set.
For each distinct kernel \(P\) choose an element \(t_P\) that
belongs to every other kernel and does not belong to \(P\).
It exists as a product of elements supplied by incomparability.
Then \(s=\sum_P t_P\) has no zero coordinate in \(Q\), hence
is a non-zero-divisor of \(F\) and a unit there. The element
\(t_P/s\) is the characteristic idempotent of the set of
coordinates with kernel \(P\). This decomposes \(F\) into domains.
If an element in one domain component is nonzero, add the
identity of all other components to obtain a non-zero-divisor
of \(F\), then use its inverse to invert the original element
in that component. The components are therefore fields.
Constancy follows from \(k\subseteq F\subseteq Q\).

The algebra \(R_F\) is finite type over a finite product of
fields, hence noetherian. It contains the factor idempotents
of \(R\) supplied by Lemma~\ref{res:PV03}. Thus any zero
coordinate of an element of \(R_F\) makes it a zero divisor
unless \(Q\) has just one factor, when the same criterion
holds in the domain. Its non-zero-divisors are precisely the
elements invertible in \(Q\). Every fraction from \(R\) can
therefore be formed from \(R_F\), and every fraction from \(R_F\)
can be formed in \(Q\). The two total rings agree.
\end{proof}

\begin{theorem}[Fixed ring over an intermediate total ring]
\label{res:PV14}
For \(F\) as in Lemma~\ref{res:PV13}, let \(G_F\) consist of the
elements of \(G\) fixing \(F\) pointwise. Then \(G_F\) is an
algebraic subgroup of \(G\) and \(Q^{G_F}=F\).
\end{theorem}

\begin{proof}
The product \(Q\otimes_F R_F\), with the operator in both
factors, contains its first factor faithfully. In fact, on
each component of \(F\) both algebras contain the component
field, and tensoring over a field preserves injections.
It is generated over \(Q\) by the constant matrix of ratios
\(Z_F=(Y\otimes1)^{-1}(1\otimes Y)\) and its inverse determinant.
The ring \(Q\) is simple as an operator ring, since it is a
field or a transitive product of fields. Lemma~\ref{res:PV08}
therefore gives
\[
 Q\otimes_F R_F \ \cong\
 Q\otimes_C C[(Z_F)_{ij},\det(Z_F)^{-1}].
\]
The algebra on the left is reduced by
Lemma~\ref{res:PV11}, since \(R_F\subseteq Q\) is reduced and
Lemma~\ref{res:PV13} identifies the base components. Thus
the constant algebra on the right is reduced.

Its \(C\)-valued evaluations induce operator maps \(R_F\to Q\)
fixing \(F\) and sending \(Y\) to \(Yg\), \(g\in GL_n(C)\).
Such a map has image exactly \(R_F\) and is a surjective
endomorphism of a noetherian ring. It is injective: the ascending
sequence of kernels of its powers stabilizes, and surjectivity
then forces the first kernel to be zero. It extends to an
automorphism of \(Q/F\) by Lemma~\ref{res:PV13}; by
Theorem~\ref{res:PV12} it belongs to \(G_F\). Conversely every
element of \(G_F\) gives such an evaluation. This realizes \(G_F\)
as the algebraic subgroup of \(G\) cut out by the relations of
\(Z_F\).

Suppose that \(u=r/s\in Q\), with \(r,s\in R\) and \(s\)
a non-zero-divisor, is fixed by \(G_F\). The element
\(1\otimes r-u\otimes s\) of \(Q\otimes_F R_F\) vanishes under
every evaluation of the constant algebra. Express this element
as a finite sum with \(C\)-linearly independent coefficients in
\(Q\). Lemma~\ref{res:PV04} implies that each of its constant
algebra coefficients is zero. Its image in \(Q\otimes_F Q\)
is therefore zero. Since \(1\otimes s\) is invertible there,
\(1\otimes u=u\otimes1\).
To see that \(u\in F\), decompose over the field factors of \(F\)
and choose in each component an \(F\)-linear map from the
corresponding component of \(Q\) to that of \(F\), sending \(1\)
to \(1\). Applying this map in the first factor to the last
tensor equality gives \(u\) in the corresponding copy of \(F\).
This proves the fixed-ring claim.
\end{proof}

\section{Separation of subgroups}

\begin{lemma}[A representation detecting a closed subgroup]
\label{res:PV15}
For a closed subgroup \(H\subseteq G\), there exist a finite
dimensional regular representation \(\rho\) of \(G\) over \(C\)
and a line \(\ell\) in its space whose stabilizer is exactly \(H\).
If \(H\) is normal, there is a finite dimensional regular
linear representation of \(G\) with kernel exactly \(H\).
\end{lemma}

\begin{proof}
The ideal of functions vanishing on \(H\) in \(C[G]\) is
finitely generated. Take a common power of the determinant in
the denominators of representatives of the generators and bound
the degrees of all their numerators. The functions with those
bounds form a finite dimensional vector space \(V\) stable
under right translation, \(\rho_0(g)f(x)=f(xg)\).
The representation is regular: translating affects each bounded
monomial by polynomial coefficients in entries of \(g\) and
\(\det(g)^{-1}\). Taking a basis of \(V\) expresses the translation
in regular matrix coefficients. Let \(W\subseteq V\) consist of
the functions vanishing on \(H\). Every \(h\in H\) preserves
\(W\). Conversely, if \(g\) preserves \(W\), each of the chosen
ideal generators \(f\in W\) satisfies \(f(g)=(\rho_0(g)f)(1)=0\),
so \(g\in H\). If \(d=\dim W\), take
\(\rho=\bigwedge^d\rho_0\) and \(\ell=\bigwedge^d W\).
For a nonzero \(w\in\ell\) the equality \(v\wedge w=0\) in
\(\bigwedge^{d+1}V\) characterizes \(v\in W\), including when
\(W=0\) or \(W=V\). Thus the stabilizer of \(\ell\) is \(H\).

Assume now \(H\) is normal. For each irreducible component
\(X\) of \(G\), let \(U_X\) be the span of the vectors
\(\rho(g)w\) with \(g\in X\). For \(h\in H\), normality
makes each such vector an eigenvector of \(\rho(h)\).
The locus where the eigenvalue is a given scalar is closed
in \(X\), and there are only finitely many eigenvalues.
Irreducibility forces one eigenvalue throughout \(X\).
Hence \(h\) acts as a scalar on \(U_X\).
Left translation by \(g\) permutes the components, so
\(\rho(g)\) permutes the finite set of distinct spaces \(U_X\).
On the direct sum of the spaces \(\operatorname{End}_C(U_X)\),
one summand for every distinct \(U_X\), act by transport and
conjugation using \(\rho(g)\). This is regular. For each possible
permutation of the summands, the set realizing it is
closed, as transport of specified subspaces is given by linear
conditions on the matrix. Its complement is a finite union of such
closed sets, so the locus is open as well. Functions regular on
these finitely many closed-and-open pieces combine into regular
functions on \(G\): the vanishing ideals of distinct pieces are
comaximal by Lemma~\ref{res:PV04}, and their intersection is zero.
For two comaximal ideals \(I,J\), write \(1=i+j\) with \(i\in I,j\in J\).
The element \(aj+bi\) then combines \(a\bmod I\) and \(b\bmod J\).
Iterating gives the finite decomposition.
On each piece, the formulas come from the regular matrices of
\(\rho(g)\) and \(\rho(g^{-1})\).
This action kills \(H\). If it kills \(g\), then \(g\)
preserves every summand and acts by a scalar on each \(U_X\):
commuting with coordinate projections makes a matrix diagonal,
and commuting with all matrix units makes the diagonal entries equal.
One of the \(U_X\) contains \(w\), so \(g\) preserves \(\ell\)
and belongs to \(H\). This is the asserted representation.
\end{proof}

\begin{lemma}[Separating a coset by a rational invariant]
\label{res:PV16}
Suppose \(H\subseteq G\) is closed and \(g\in G\setminus H\).
Then \(g\) acts nontrivially on \(Q^H\).
\end{lemma}

\begin{proof}
First consider the coordinate algebra \(C[G]\) and its total
ring of fractions. Take \(\rho,\ell=Cw\) as in
Lemma~\ref{res:PV15}. Choose a linear functional \(a\)
nonzero on \(w\), on \(\rho(g)w\), and on at least one vector
\(\rho(x)w\) from each irreducible component of \(G\).
This is possible over the infinite field \(C\) by avoiding
finitely many proper hyperplanes. Thus \(a(\rho(x)w)\)
is a non-zero-divisor in \(C[G]\). Since the lines of
\(w\) and \(\rho(g)w\) are distinct, choose a linear functional
\(b\) that takes different values on their vectors normalized
by \(a=1\). The rational function
\[
 f(x)=\frac{b(\rho(x)w)}{a(\rho(x)w)}
\]
is invariant under right translations by \(H\), because those
translations multiply numerator and denominator by the same
nonzero scalar. It has distinct values at \(1\) and \(g\).
Consequently right translation by \(g\) does not fix \(f\).

We transport this separation from \(G\) to \(Q\). Choose a
maximal ideal in \(R\) and let \(E\) be its residue field.
It is a finite extension of \(k\) by Lemma~\ref{res:PV04}.
Base-changing the first factor of Theorem~\ref{res:PV12}
along \(R\to E\) gives
\[
 E\otimes_k R\cong E\otimes_C C[G].
\]
Here an element \(g\in G\) acts on the first algebra by its
action on \(R\) and on the second by right translation,
because \(Y_1^{-1}Y_2\) changes to \(Y_1^{-1}Y_2g\).
The total ring of fractions on the left is \(E\otimes_k Q\).
Indeed this tensor is a product of fields, since
\(E/k\) is finite separable and each factor of \(Q\) is a field.
It is a localization of \(E\otimes_k R\), at non-zero-divisors
from \(R\). Any non-zero-divisor of \(E\otimes_k R\) remains
a non-zero-divisor after localization and hence is a unit
in the resulting product of fields. On the right, the
denominator of \(f\) stays a non-zero-divisor after extension
from \(C\) to \(E\), by exactness of tensoring over a field.
Thus \(f\) gives an \(H\)-invariant element of \(E\otimes_k Q\)
moved by \(g\); the distinct evaluations at \(1,g\) persist
over \(E\). Since invariants of \(E\otimes_k Q\) are
\(E\otimes_k Q^H\), as one sees by expanding in a \(k\)-basis
of \(E\), the action of \(g\) on \(Q^H\) is nontrivial.
\end{proof}

\begin{theorem}[Differential and difference Galois correspondence]
\label{res:PV17}
The maps
\[
 H\longmapsto Q^H,\qquad
 F\longmapsto G_F
\]
are inverse inclusion-reversing bijections between closed
subgroups of \(G\) and intermediate operator subrings \(F\)
in which every non-zero-divisor is a unit. In the differential
case the latter are precisely the intermediate differential
fields. In particular \(Q^G=k\).
\end{theorem}

\begin{proof}
For any subgroup \(H\), its fixed subring is stable under the
operator and contains \(k\). If an element \(u\in Q^H\) vanishes
in some component of \(Q\), the characteristic idempotent
of its zero components is nonzero, belongs to \(Q^H\), and
annihilates \(u\). Indeed every element of \(H\) permutes
the field factors and preserves this zero-component set.
Thus a non-zero-divisor of \(Q^H\) has no zero component;
its inverse in \(Q\) is fixed by \(H\). This proves that
\(Q^H\) is an intermediate ring of the indicated kind.
For such a ring \(F\), Theorem~\ref{res:PV14} gives \(Q^{G_F}=F\).
For a closed subgroup \(H\), each element of \(G\setminus H\)
is excluded from \(G_{Q^H}\) by Lemma~\ref{res:PV16}, and
\(H\subseteq G_{Q^H}\) by definition. The inclusion reversal
is immediate from the two definitions. Taking \(F=k\) gives
the full fixed-ring equality.
\end{proof}

\section{Normal subgroups and Picard--Vessiot quotients}

\begin{lemma}[Finite semilinear descent of a matrix]
\label{res:PV18}
Over a field \(k\) of arbitrary characteristic, suppose \(E/k\) is
finite Galois, with group \(\Gamma\), and \(a_\gamma\in GL_m(E)\) obeys
\(a_{\gamma\eta}=a_\gamma\gamma(a_\eta)\).
Then for some \(B\in GL_m(E)\) one has
\(\gamma(B)=B a_\gamma\) for all \(\gamma\in\Gamma\).
\end{lemma}

\begin{proof}
On row vectors in \(E^m\) put \(T_\gamma(v)=\gamma(v)a_\gamma^{-1}\).
The cocycle identity makes this an action by semilinear maps.
The sums \(P(v)=\sum_\gamma T_\gamma(v)\) are invariant.
They span \(E^m\) over \(E\): if a column \(c\) annihilates all
of them, substitute \(v=x e_i\) for each coordinate row \(e_i\).
Then \(\sum_\gamma \gamma(x)e_i a_\gamma^{-1}c=0\) for all \(x\).
Distinct field automorphisms are linearly independent as functions
on \(E\) by~\cite[Lemma 16.6, PDF p.~125]{mg-d07}.
Thus the coefficient of the identity automorphism gives
\(e_i c=0\) for all \(i\).
Taking \(m\) independent invariant rows gives \(B\) as required.
\end{proof}

\begin{theorem}[The normal subgroup criterion and quotient]
\label{res:PV19}
An intermediate total ring \(F=Q^H\) of Theorem~\ref{res:PV17}
is stable under the action of \(G\) if and only if \(H\)
is normal. In that case \(F\) is the total extension of a
Picard--Vessiot ring for a system over \(k\), and restriction
is a regular morphism in the matrix realizations and gives an
isomorphism of groups of \(C\)-points
\[
 G/H \ \cong\ \operatorname{Aut}_{\mathrm{op}}(F/k).
\]
Conversely, any intermediate total ring that is a
Picard--Vessiot total extension over \(k\) is stable under \(G\).
\end{theorem}

\begin{proof}
The correspondence of Theorem~\ref{res:PV17} and
\(g(Q^H)=Q^{gHg^{-1}}\) give the stability criterion.
If \(H\) is normal, take a regular representation
\(\rho:G\to GL_m(C)\) with kernel \(H\) from
Lemma~\ref{res:PV15}. There is a finite Galois extension \(E/k\)
over which the torsor has an \(E\)-point: take a finite residue
extension as in Lemma~\ref{res:PV16} and then a splitting
field of its separable defining polynomial. If \(U\in GL_n(E)\)
denotes the image of \(Y\) at this point, the torsor coordinates
on \(E\otimes_k R\) are \(Z=U^{-1}Y\). Writing
\(\gamma(U)=U c_\gamma\), the torsor formula puts
\(c_\gamma\in G(E)\) and gives
\(c_{\gamma\eta}=c_\gamma\gamma(c_\eta)\). Here \(G(E)\)
means the \(E\)-valued solutions of the relations for \(G\).
The multiplication law and the regular representation extend
to points valued in any \(C\)-algebra. Indeed their polynomial identities on
\(C\)-points are identities in the coordinate algebras,
by Lemma~\ref{res:PV04}, also for products. For the latter
assertion, successive evaluation in each factor of a
tensor product of coordinate algebras detects any nonzero
element by independence of coefficients.

By Lemma~\ref{res:PV18} take \(B\) with
\(\gamma(B)=B\rho(c_\gamma)\). Then
\[
 Y_H=B\rho(Z)\in GL_m(E\otimes_k R)
\]
is invariant under \(\Gamma\), since
\(\gamma(Z)=c_\gamma^{-1} Z\).
Expanding tensor elements in a \(k\)-basis of \(R\) proves
that \((E\otimes_k R)^\Gamma=R\), so \(Y_H\) and its inverse
are in \(R\). The action of an element \(g\in G\) on this
matrix is \(Y_H\mapsto Y_H\rho(g)\). Its entries are therefore
fixed by \(H\). The coefficient matrix
\(\partial(Y_H)Y_H^{-1}\) or \(\sigma(Y_H)Y_H^{-1}\) is
fixed by all of \(G\), as the right factors \(\rho(g)\)
are constant. By Theorem~\ref{res:PV17} it lies in \(M_m(k)\);
in the difference case it is invertible there, by applying
the same argument to its inverse. Therefore \(Y_H\) is a
fundamental matrix for a system over \(k\).

Its ring \(R_H=k[(Y_H)_{ij},\det(Y_H)^{-1}]\) is
Picard--Vessiot by Theorem~\ref{res:PV09}, applied to \(Q\).
Each non-zero-divisor of \(R_H\) is invertible in \(Q\).
For the difference case, use the distinct projection kernels
from \(R_H\) to the field factors of \(Q\). They are
transitively permuted and hence incomparable, as in the proof
of Lemma~\ref{res:PV13}. If an element vanishes in one
projection, it annihilates the nonzero product belonging to
all other distinct kernels and outside that kernel. Thus
every non-zero-divisor has nonzero image in each factor.
It follows that \(Q(R_H)\) is an intermediate total ring in
\(Q^H\). An automorphism fixes \(Q(R_H)\) precisely if it
fixes \(Y_H\), or equivalently if it belongs to \(H\).
Theorem~\ref{res:PV17} now yields \(Q(R_H)=F\).

In the matrix realizations restriction is \(g\mapsto\rho(g)\),
so it is regular and its kernel is \(H\).
To prove surjectivity, let \(\tau\) be an operator
automorphism of \(F/k\). Form \(Q\otimes_F R_F\) with the
map \(F\to Q\) on the first side twisted by \(\tau\).
It is nonzero and contains \(Q\) by the field-component
argument of Theorem~\ref{res:PV14}. It is generated over
\(Q\) by the same constant ratios of fundamental matrices.
Lemma~\ref{res:PV08} and a \(C\)-point of their constant
algebra give a map \(R_F\to Q\) extending \(\tau\) and sending
\(Y\) to \(Yg\) for some \(g\in GL_n(C)\). Its image is
\(R_F\); it is an isomorphism by the noetherian kernel argument
in Theorem~\ref{res:PV14}. It extends to \(Q\), yielding
the desired lift in \(G\). Finally, if an intermediate ring is
given as a Picard--Vessiot total extension, a fundamental
matrix for it is carried by every \(g\in G\) to another
fundamental matrix for the same system. Its constant ratio
lies in \(GL_m(C)\). The ring of entries and its total ring
are preserved, proving the converse.
\end{proof}

\begin{corollary}[The algebraic part of a differential extension]
\label{res:PV20}
In the differential case, let \(G^0\) be the connected component
of the identity in the Zariski topology of \(G\). Then \(Q^{G^0}\)
is the relative algebraic closure of \(k\) in \(Q\),
and it is a finite Galois extension of \(k\) with group \(G/G^0\).
\end{corollary}

\begin{proof}
The component \(G^0\) is a closed normal subgroup.
Indeed inversion and conjugation preserve the identity component,
and for \(g\in G^0\) the connected set \(gG^0\) meets \(G^0\) at
\(g\), hence lies in \(G^0\). The components are its translates, and
there are finitely many since each is a union of irreducible
components of a noetherian space.
For \(u\in Q\) algebraic over \(k\), its orbit is finite because
all its translates are roots of its minimal polynomial.
Its stabilizer is the closed subgroup \(G_{k(u)}\) by
Theorem~\ref{res:PV14}. A closed subgroup of finite index is also
open because its complement is a finite union of its other cosets.
Thus the stabilizer contains \(G^0\), proving \(u\in Q^{G^0}\).

Conversely, every \(u\in Q^{G^0}\) has finite \(G\)-orbit of
cardinality at most \(d=|G/G^0|\). The polynomial whose roots are its
distinct translates has coefficients in \(Q^G=k\) by
Theorem~\ref{res:PV17}. It splits in \(Q^{G^0}\) and is separable.
Thus each \(u\) is algebraic of degree at most \(d\), and its minimal
polynomial splits in \(Q^{G^0}\). Any finite subextension is
separable and has a primitive element by
\cite[Section 14.31, PDF p.~113]{mg-d07}, so also has degree at most
\(d\). A finite subextension of largest degree must contain
every \(u\), since adjoining any further element would increase the
degree. Hence \(Q^{G^0}/k\) is finite, normal, and separable. The
automorphisms in Theorem~\ref{res:PV19} are all its field
automorphisms, since a derivation extends uniquely to a finite
separable field extension. Uniqueness follows by differentiating
the minimal polynomial of each element, whose derivative at that
element is nonzero.
\end{proof}

\section{Two rank-one calculations}

\begin{corollary}[The rank-one multiplicative group]
\label{res:PV21}
Consider \(\partial y=a y\), \(a\in k\), or
\(\sigma(y)=a y\), \(a\in k^\times\). The Galois group is
either \(C^\times\) or the group \(\mu_m(C)\) of \(m\)-th roots
of unity. The finite case occurs exactly when there is some
positive \(m\) and \(b\in k^\times\) with
\[
 \partial b=ma b \quad\text{or}\quad \sigma(b)=a^m b,
\]
respectively. The least such \(m\) is the order of the group.
\end{corollary}

\begin{proof}
Every element \(g\in G\) acts on a fundamental unit \(y\) by
\(y\mapsto gy\). By Theorem~\ref{res:PV12}, \(G\) is a
closed subgroup of \(C^\times\), so it is finite unless it is
all of \(C^\times\), since a nonzero Laurent polynomial has
finitely many roots. A finite subgroup of order \(m\) lies
in \(\mu_m(C)\) by the group order theorem, hence equals it.
The element \(y^m\) is fixed exactly when \(G\subseteq\mu_m(C)\),
so by Theorem~\ref{res:PV17} this condition is equivalent to
\(y^m\in k^\times\). Then \(b=y^m\) satisfies the required
equation. Conversely, \(y^m/b\) is constant for any such
\(b\) and belongs to \(C^\times\), hence \(y^m\in k^\times\).
This proves both the criterion and its minimality assertion.
\end{proof}

\begin{example}[A disconnected difference Galois group]
\label{res:PV22}
The extension of Example~\ref{res:PV10} has Galois group
\(\{1,-1\}\), acting by \(u\mapsto\pm u\), and its full fixed
ring is \(k\). For the differential field \(C(z)\) with
\(\partial z=1\), the equation \(\partial y=y\) has Galois
group \(C^\times\).
\end{example}

\begin{proof}
In Example~\ref{res:PV10}, the relation \(u^2=1\) forces the
multiplier to have square \(1\); both choices give
automorphisms. The nonidentity automorphism interchanges the
two factors of \(k\times k\) and fixes its diagonal.
In the differential example the constants are \(C\), since
a rational function of derivative zero in characteristic zero
cannot have a finite pole and then must be a constant polynomial.
For any nonzero relatively prime polynomials \(p,q\), the
logarithmic derivative of \(p/q\) is
\((p'q-pq')/(pq)\), whose numerator has smaller degree than
its denominator whenever it is nonzero. It cannot equal
a positive integer \(m\) in characteristic zero, so
Corollary~\ref{res:PV21} gives the claim.
\end{proof}

\paragraph{Sources.}
For differential Galois groups and the correspondence see
van der Put--Singer~\cite[Chapter 1, Section 1.4]{vps-differential}
and Singer~\cite[Sections 1.3 and 1.5]{singer}. For the total
ring and difference correspondence see van der Put--Singer
\cite[Chapter 1, Sections 1.2--1.3]{vps-difference} and
Roques~\cite[Sections 5.2--8]{roques}.
\chapter{Galois Cohomology and Field Arithmetic}
\label{ch:galois-cohomology}

Throughout, \(\operatorname{Gal}(L/k)\) for a possibly infinite Galois
extension has the Krull topology. Fix a separable closure \(k^{\mathrm s}\)
and write \(G_k=\operatorname{Gal}(k^{\mathrm s}/k)\). Superscript
\({}^{+}\) denotes an additive group. The coefficient groups and matrix
groups attached to a field in this chapter have the discrete topology.
We use the continuous inhomogeneous complex of Result GC01 in
\emph{Topics in Group Theory} \cite[PDF p.~24]{mg-groups}. In degree one,
with multiplicative notation, \(c(gh)=c(g)\,g(c(h))\); for abelian
coefficients the coboundary attached to \(b\) is \(g\mapsto b^{-1}g(b)\).
For a possibly nonabelian coefficient group the equivalence is
\[
 c'(g)=b^{-1}c(g)\,g(b).
\]
The resulting \(H^1\) is a pointed set in this case, with the identity
cocycle as base point.

\section{Finite layers and Hilbert's Theorem 90}

\begin{proposition}[Finite Galois layers and fixed coefficients]
\label{res:GH01}
Let \(L/k\) be Galois and put \(G=\operatorname{Gal}(L/k)\).
For each finite Galois intermediate field \(E\) put
\(U_E=\operatorname{Gal}(L/E)\) and
\(G_E=\operatorname{Gal}(E/k)\). For a discrete continuous abelian
\(G\)-module \(M\) there are identifications
\[
 H^q_{\mathrm{cts}}(G,M)
   \ \cong\ \varinjlim_E H^q(G_E,M^{U_E})\qquad(q\geq 0).
\]
They use inflation along restriction of field automorphisms and inclusion
of fixed coefficients. For a discrete group \(B\) with continuous
\(G\)-action, the analogous identification holds for the pointed sets
\(H^1_{\mathrm{cts}}(G,B)\).
\end{proposition}
\begin{proof}
Results F04, T02, T06 and C02 of \emph{Infinite Galois Theory}
\cite[PDF pp.~5, 9--10, 12]{mg-aft} give the directed finite layers,
the open normal basis \(U_E\), the quotient \(G/U_E=G_E\), and
\(L^{U_E}=E\). The abelian assertion is GC02 of \emph{Topics in Group
Theory} \cite[PDF p.~24]{mg-groups} applied to this cofinal basis.

For a nonabelian cocycle, compactness gives a finite image, and
Result PF03(4) of \emph{Topics in Group Theory}
\cite[PDF p.~16]{mg-groups} factors its underlying continuous map
through a finite quotient. Enlarge the layer \(E\) so that \(U_E\)
fixes all image elements, by intersecting their open stabilizers
with the first quotient kernel and choosing a basis subgroup.
Then the cocycle identity descends by surjectivity of restriction.
If two cocycles are related by some \(b\in B\), enlarge the common
layer once more so that \(U_E\) fixes \(b\). Their equivalence now
holds at that stage. These two finite-layer assertions give
surjectivity and injectivity of the pointed-set comparison.
\end{proof}

\begin{lemma}[A trace-one element]
\label{res:GH02}
If \(E/k\) is finite Galois, there exists \(t\in E\) with
\(\sum_{\sigma\in\operatorname{Gal}(E/k)}\sigma(t)=1\).
\end{lemma}
\begin{proof}
The independence lemma for distinct field embeddings in D07,
Lemma 16.6 \cite[PDF p.~125]{mg-d07}, makes the sum of the
automorphisms in the finite Galois group a nonzero function.
It is \(k\)-linear and its values are fixed by the Galois group, hence
are in \(k\) by the finite fixed-field theorem (Result F01 of
\emph{Infinite Galois Theory} \cite[PDF p.~4]{mg-aft}).
Divide an element of nonzero trace by its trace.
\end{proof}

\begin{theorem}[Matrix Hilbert 90]
\label{res:GH03}
For \(n\geq 1\) and a Galois extension \(L/k\), with entrywise action,
\[
 H^1_{\mathrm{cts}}\bigl(\operatorname{Gal}(L/k),
                         \operatorname{GL}_n(L)\bigr)=\{*\}.
\]
In particular \(H^1_{\mathrm{cts}}(\operatorname{Gal}(L/k),L^\times)=0\).
\end{theorem}
\begin{proof}
For a finite extension \(E/k\) and a cocycle with values in
\(\operatorname{GL}_n(E)\), Lemma~\ref{res:PV18} gives an invertible
matrix \(b\) with \(A_\sigma=b^{-1}\sigma(b)\) in any characteristic.
The stated equivalence applied to the identity cocycle shows
that this cocycle represents the base point.
In the infinite case \(\operatorname{GL}_n(L)^{U_E}=
\operatorname{GL}_n(E)\): the fixed entries lie in \(E\), and a
determinant nonzero in \(L\) is nonzero in \(E\).
Proposition~\ref{res:GH01} now reduces the assertion to the finite case.
\end{proof}

\begin{proposition}[Evaluation at a generator in degree one]
\label{res:GH04}
For \(C=\langle s\rangle\) of finite order \(d\) acting on an abelian
group \(A\), evaluation at \(s\) gives an isomorphism
\[
 H^1(C,A)\ \cong\ \ker(N:A\to A)/(s-1)A,\qquad
 N=1+s+\cdots+s^{d-1}.
\]
\end{proposition}
\begin{proof}
If \(d=1\), both sides are zero. Otherwise a cocycle with value \(a\)
at \(s\) has value \(\sum_{j=0}^{i-1}s^j a\) at \(s^i\) for
\(0\leq i<d\). At \(i=d\) the cocycle identity gives \(Na=0\);
conversely this condition makes the sum rule consistent under
reduction modulo \(d\) and implies the cocycle identity. A
coboundary has value \((s-1)b\) at \(s\), so quotienting gives
the stated isomorphism.
\end{proof}

Thus the older multiplicative and additive forms of cyclic
Hilbert 90 are statements that these particular quotients vanish.
D07 proves the norm-one and trace-zero formulas in Results
\texttt{K07.20.hilbert90\_multiplicative} and
\texttt{K07.20.hilbert90\_additive} \cite[PDF pp.~171--172]{mg-d07}.

\section{Kummer theory}

\begin{proposition}[Exact coefficient sequences for continuous cochains]
\label{res:GH05}
An exact sequence \(0\to A\to B\to C\to0\) of discrete abelian
continuous \(G\)-modules, for \(G\) profinite, gives the long exact
sequence in continuous cohomology. The connecting arrow sends
\([c]\) to \([d\widetilde c]\), for a continuous cochain lift
\(\widetilde c\) to \(B\).
\end{proposition}
\begin{proof}
The induced cochain map to \(C\) is surjective: a continuous map from
the compact space \(G^n\) to a discrete group takes finitely many
values, and the choice of one lift for each value gives a continuous
lift. The two other exactness assertions hold pointwise, with the
induced topology on each subgroup still discrete. This gives a
short exact sequence of cochain complexes. Its connecting arrow
is obtained by lifting a cocycle and applying the differential,
as in \emph{Algebra Core}, Theorem 10.3
\cite[PDF p.~51]{mg-algebra-core}.
\end{proof}

\begin{theorem}[Kummer isomorphism with its Galois action]
\label{res:GH06}
Let \(m\geq 1\) and assume that \(m\) is invertible in \(k\). Give
\(\mu_m\subset(k^{\mathrm s})^\times\) its natural \(G_k\)-action.
Then there is an isomorphism
\[
 k^\times/k^{\times m}
   \ \xrightarrow{\ \sim\ }\ H^1_{\mathrm{cts}}(G_k,\mu_m),
 \qquad
 a\longmapsto [\,g\mapsto g(\alpha)/\alpha\,],\quad \alpha^m=a.
\]
If \(\mu_m\subset k^\times\), this identifies the quotient with
\(\operatorname{Hom}_{\mathrm{cts}}(G_k,\mu_m)\).
\end{theorem}
\begin{proof}
The \(m\)-th power on \((k^{\mathrm s})^\times\) is surjective.
Indeed \(X^m-u\) is separable for \(u\ne0\), by its derivative, and
splits over the separably closed \(k^{\mathrm s}\). Its kernel on
units is \(\mu_m\). The maps are equivariant, and all coefficients
are continuous by the stabilizers in the Krull topology.
Invariants in \((k^{\mathrm s})^\times\) are \(k^\times\), by
Result C02 of \emph{Infinite Galois Theory}
\cite[PDF p.~12]{mg-aft}. Proposition~\ref{res:GH05} gives the connecting
map of the short exact sequence, and Theorem~\ref{res:GH03} makes it
surjective onto \(H^1_{\mathrm{cts}}(G_k,\mu_m)\). Its formula follows
from the coboundary of \(\alpha\). Changing \(\alpha\) by an element of
\(\mu_m\) multiplies the cocycle by that element's coboundary. For
trivial action the cocycle identity is the homomorphism condition and
every coboundary is the identity.
\end{proof}

\begin{theorem}[Finite Kummer duality]
\label{res:GH07}
Assume that \(m\) is invertible in \(k\) and
\(\mu_m\subset k^\times\). Inside \(k^{\mathrm s}\), the finite Galois
extensions with abelian group of exponent dividing \(m\)
correspond to the finite subgroups \(D\subset k^\times/k^{\times m}\).
The field \(L_D\) is obtained by adjoining \(m\)-th roots of
representatives of \(D\). If \(L\) is one of the extensions, its
subgroup is
\[
 D_L=\bigl(k^\times\cap (L^\times)^m\bigr)/k^{\times m}.
\]
The pairing \((g,[a])\mapsto g(\alpha)/\alpha\), with \(\alpha^m=a\),
identifies either of \(\operatorname{Gal}(L_D/k)\) and \(D\) with the
group of homomorphisms from the other to \(\mu_m\).
\end{theorem}
\begin{proof}
Any finite subgroup of the units of a field is cyclic. Its exponent
\(e\) is realized as the order of an element: choose the prime-power
parts of elements having maximal such parts and multiply them.
Then the polynomial \(X^e-1\) bounds the order of the subgroup by \(e\).
In particular \(\mu_m\) is cyclic of order \(m\), since \(X^m-1\) is
separable and split in \(k^{\mathrm s}\).

For a finite abelian group \(A\) of exponent dividing \(m\), its
\(\mu_m\)-valued characters separate elements and number \(|A|\).
Indeed, a character on a subgroup \(H\) extends across an element
\(x\) of order \(d\) modulo \(H\). Since \(d\mid m\),
\(\chi(x^d)^{m/d}=1\); choose \(\zeta\in\mu_m\) with
\(\zeta^d=\chi(x^d)\). Then
\(\chi(hx^i)=\chi(h)\zeta^i\) gives the extension. Repeating proves
surjectivity of restriction to any subgroup. This extends a faithful
character of each cyclic subgroup. Counting the characters by
restriction to a nontrivial cyclic subgroup and using the quotient
for the kernel gives the count by induction on \(|A|\).

Choose finitely many generators of \(D\). The field generated by
each chosen root is finite cyclic Galois by the single-radical
result \texttt{K07.15.kummer\_forward} in D07
\cite[PDF p.~135]{mg-d07}. By Result F04 of
\emph{Infinite Galois Theory} \cite[PDF p.~5]{mg-aft},
their compositum \(L_D/k\) is finite Galois.
If \([a]\) is a product of powers of the generator classes, an
\(m\)-th root of \(a\) in \(L_D\) is the corresponding product of
powers of the chosen roots times a scalar in \(k^\times\).
On these generators the ratios give an injective homomorphism
of \(\operatorname{Gal}(L_D/k)\) into a product of \(\mu_m\)'s.
The pairing is independent of roots and representatives and is
multiplicative in both variables. Its right kernel is trivial by
the fixed-field theorem: if each \(g\) fixes \(\alpha\), then
\(\alpha\in k\), so \([a]=1\). The left kernel is trivial because
the chosen roots generate \(L_D\). The character counts just proved
now make both maps isomorphisms.

For a finite Galois \(L/k\) with abelian group \(A\) of exponent
dividing \(m\), each character \(\chi:A\to\mu_m\) is a cocycle on
\(L^\times\). Theorem~\ref{res:GH03} gives \(\chi(g)=g(b)/b\), with
\(b\in L^\times\) and \(b^m\in k^\times\). Choosing such a \(b\) for
each character generates \(L\): the subgroup fixing them all is
trivial because the characters separate, and the finite Galois
correspondence applies (Result F01 of \emph{Infinite Galois Theory}
\cite[PDF p.~4]{mg-aft}). Finally \([a]\mapsto(g\mapsto g(\alpha)/\alpha)\)
injects \(D_L\) into the character group, by the fixed-field theorem,
and the chosen elements show surjectivity. For \(L=L_D\), inclusion
\(D\subseteq D_L\) and the equal sizes from perfectness give equality.
\end{proof}

\section{Artin--Schreier theory}

\begin{proposition}[Trace contraction for semilinear coefficients]
\label{res:GH08}
Let \(E/k\) be finite Galois with group \(H\), and let \(V\) be an
\(E\)-vector space on which \(H\) acts semilinearly, so
\(\sigma(av)=\sigma(a)\sigma(v)\). Its additive group has
\(H^q(H,V^+)=0\) for all \(q>0\). For any Galois extension \(L/k\),
\[
 H^q_{\mathrm{cts}}\bigl(\operatorname{Gal}(L/k),L^+\bigr)=0
 \qquad(q>0).
\]
\end{proposition}
\begin{proof}
Use equivariant homogeneous cochains \(f:H^{q+1}\to V^+\),
with differential the alternating sum of omissions.
The change to inhomogeneous cochains evaluates \(f\) at
\((1,g_1,g_1g_2,\ldots,g_1\cdots g_q)\); its inverse takes
\(F\) to \(g_0 F(g_0^{-1}g_1,\ldots,g_{q-1}^{-1}g_q)\).
Deleting the first variable yields the action term and deleting
an interior variable multiplies adjacent increments. These formulas
intertwine the two differentials.
Choose trace-one \(t\) by Lemma~\ref{res:GH02}. For \(q\geq1\) put
\[
 (s f)(g_0,\ldots,g_{q-1})
   =\sum_{h\in H}h(t)\,f(h,g_0,\ldots,g_{q-1}).
\]
The reindexing \(h=\sigma h'\) and semilinearity prove
equivariance. In \(s d f\), omitting \(h\) gives
\(\sum_h h(t)f=f\); each of the remaining terms cancels its
counterpart in \(d s f\) by the shift of signs. Hence \(d s+s d=1\)
on positive degrees. For \(L/k\) the fixed coefficients at stage
\(E\) are \(E^+\), so use Proposition~\ref{res:GH01}.
\end{proof}

\begin{theorem}[The Artin--Schreier connecting isomorphism]
\label{res:GH09}
If \(\operatorname{char}k=p>0\) and \(\wp(x)=x^p-x\), then
\[
 k/\wp(k)
   \ \cong\ H^1_{\mathrm{cts}}(G_k,\mathbb F_p)
   =\operatorname{Hom}_{\mathrm{cts}}(G_k,\mathbb F_p),
 \qquad
 [a]\longmapsto \bigl(g\mapsto g(\alpha)-\alpha\bigr)
 \quad(\wp(\alpha)=a).
\]
The coefficient action on \(\mathbb F_p\) is trivial.
\end{theorem}
\begin{proof}
The polynomial \(X^p-X-u\) is separable, since its derivative is
\(-1\), so \(\wp:k^{\mathrm s}\to k^{\mathrm s}\) is surjective.
Its kernel is exactly \(\mathbb F_p\), which already gives \(p\)
roots. The resulting sequence of discrete coefficients is exact.
Apply Proposition~\ref{res:GH05}, the invariants formula
(Result C02 of \emph{Infinite Galois Theory}
\cite[PDF p.~12]{mg-aft}), and Proposition~\ref{res:GH08} in
degree one. The connecting cocycle is the difference displayed.
\end{proof}

For a nonzero class its character has kernel
\(\operatorname{Gal}(k^{\mathrm s}/k(\alpha))\), since the value at
\(g\) is zero exactly when \(g\) fixes \(\alpha\). This realizes
the cyclic degree-\(p\) extensions of D07,
\texttt{K07.20.artin\_schreier} \cite[PDF p.~172]{mg-d07},
as the lines in the character space: changing the isomorphism from
the cyclic quotient to \(\mathbb F_p^+\) rescales the character.

\begin{theorem}[Vanishing for \(p\)-primary coefficients in characteristic \(p\)]
\label{res:GH10}
In characteristic \(p>0\), every discrete continuous \(G_k\)-module
\(M\) with \(p\)-primary torsion underlying group satisfies
\[
 H^q_{\mathrm{cts}}(G_k,M)=0\qquad(q\geq2).
\]
\end{theorem}
\begin{proof}
First suppose \(pM=0\), and put
\(W=k^{\mathrm s}\otimes_{\mathbb F_p}M\), with diagonal action and
the discrete topology. It is continuous because each tensor is a
finite sum of elementary tensors, and each factor has an open
stabilizer. The Artin--Schreier sequence tensored over \(\mathbb F_p\)
with \(M\) gives
\[
 0\longrightarrow M\xrightarrow{m\mapsto 1\otimes m}W
 \xrightarrow{\wp\otimes 1}W\longrightarrow0.
\]
Exactness follows by choosing a vector-space basis of \(M\), which
turns the tensor sequence into a direct sum of copies of the
Artin--Schreier sequence from Theorem~\ref{res:GH09}. The maps in the
display are equivariant because \(\wp\) commutes with field
automorphisms.

For any continuous \(q\)-cocycle with values in \(W\), finitely
many field elements and vectors occur in expressions for its
finitely many values. A vector in the discrete \(M\) has finite
orbit by compactness of \(G_k\). Thus there is a finite-dimensional
stable \(\mathbb F_p\)-subspace \(M_0\) containing those vectors.
Choose a finite Galois layer \(E\) containing the field elements,
large enough that \(U_E\) fixes \(M_0\) and the cocycle factors
through \(G_E^q\). For the latter use the cochain factorization in
the proof of Proposition~\ref{res:GH01}. For the former use the
intersection of the stabilizers of a finite basis. The cocycle
then takes values in \(E\otimes_{\mathbb F_p} M_0\), on which
\(G_E\) acts semilinearly. This group embeds in \(W\) since tensors
over a field preserve injections, and restriction onto \(G_E\)
is surjective. The cocycle identity holds at that stage. For
\(q>0\), Proposition~\ref{res:GH08} kills the class at that stage,
so \(H^q_{\mathrm{cts}}(G_k,W)=0\). The sequence above and
Proposition~\ref{res:GH05} prove the assertion for \(pM=0\).

For \(M\) killed by \(p^r\), use the exact sequence
\(0\to M[p]\to M\xrightarrow{p}pM\to0\) and induction on \(r\).
For general \(p\)-primary torsion \(M\), each continuous cocycle
has finite image and takes values in \(M[p^r]\) for some \(r\).
Its coboundary primitive found in that submodule also works in \(M\).
\end{proof}

\section{Cyclic \texorpdfstring{\(H^2\)}{H2} and arithmetic calculations}

\begin{proposition}[Second cohomology of a finite cyclic group]
\label{res:GH11}
Let \(C=\langle s\rangle\) have order \(d\geq1\) and \(A\) be an
abelian \(C\)-module. For \(N=1+s+\cdots+s^{d-1}\), one has
\[
 H^2(C,A)\cong A^C/N(A).
\]
In particular for a finite cyclic field extension \(E/k\),
\[
 H^2\bigl(\operatorname{Gal}(E/k),E^\times\bigr)
     \cong k^\times/N_{E/k}(E^\times).
\]
\end{proposition}
\begin{proof}
The order-one case is immediate from the one-element group complex.
For \(d\geq2\) use additive notation. Every class has a normalized
cocycle \(c\), with \(c(1,g)=c(g,1)=0\). In fact the cocycle identity
implies \(c(1,g)=c(1,1)\) and \(c(g,1)=g c(1,1)\); subtract the
coboundary of the constant one-cochain \(c(1,1)\).
Define a group on \(A\times C\) by
\((a,g)(b,h)=(a+g b+c(g,h),gh)\).
Associativity is the cocycle identity, the normalized values give
the identity \((0,1)\), and inverse elements are obtained by first
solving the equation in \(C\) and then in \(A\).
The element \(u=(0,s)\) conjugates \(A\) by \(s\), and
\(u^d=(a,1)\) for some \(a\in A^C\), because \(u\) commutes with
\(u^d\). A different lift \(b u\), for \(b\in A\), replaces
\(a\) by \(a+N b\) after moving the \(d\) copies of \(b\) to the
left. Replacing \(c\) by \(c+df\) induces the isomorphism
\((b,g)\mapsto (b+f(g),g)\) from the new group to the old one, so the
class of \(a\) is well defined.

Use the section \(s^i\mapsto u^i\) for \(0\leq i<d\).
Changing a section changes its factor set by a coboundary:
if each section element is multiplied on the left by \(f(g)\in A\),
its factor set at \((g,h)\) changes by
\(f(g)+g f(h)-f(gh)\). The factor set of this section has value \(0\)
at \((s^i,s^j)\) when \(i+j<d\) and \(a\) otherwise.
Conversely, for every \(a\in A^C\) this carry rule defines a
cocycle: adding \(i,j,\ell\) and reducing modulo \(d\) in either
order gives the same total carry because \(a\) is invariant.
Its \(d\)-th power computation gives \(a\), and its cohomology
class depends additively on \(a\). If two such cocycles have
parameters differing by \(N b\), the section defined by the
new lift \(b u\) changes one into the other. These checks give the
isomorphism. The field formula follows by the finite fixed-field
theorem and the multiplicative norm.
\end{proof}

\begin{example}[Finite fields]
\label{res:GH12}
For \(k=\mathbb F_q\) of characteristic \(p\) and \(p\nmid m\),
\(H^1_{\mathrm{cts}}(G_k,\mu_m)\) is cyclic of order
\(\gcd(m,q-1)\). The Artin--Schreier quotient \(k/\wp(k)\)
is identified with \(\mathbb F_p\) by the trace, and the
character of \(a\) takes the trace value at the arithmetic
Frobenius \(x\mapsto x^q\). Also
\[
 H^2_{\mathrm{cts}}\bigl(G_k,(k^{\mathrm s})^\times\bigr)=0.
\]
\end{example}
\begin{proof}
Apply Theorem~\ref{res:GH06} and cyclicity of the units, proved in
Theorem~\ref{res:GH07}. Write \(q=p^r\). The trace
\(\operatorname{Tr}_{\mathbb F_q/\mathbb F_p}\) is
\(\sum_{i=0}^{r-1}a^{p^i}\), by the Frobenius description of
finite Galois layers (Results E01--E04 of \emph{Infinite Galois
Theory} \cite[PDF pp.~18--19]{mg-aft}), and induces the claimed
quotient by \texttt{K07.20.finite\_artin\_schreier} of D07
\cite[PDF p.~173]{mg-d07}. For a root \(\wp(\alpha)=a\),
the equality \(\alpha^q-\alpha=\sum_{i=0}^{r-1}a^{p^i}\) follows
by raising the defining equation to \(p^i\) and summing. This gives
the stated evaluation in Theorem~\ref{res:GH09}.
If \(E=\mathbb F_{q^d}\), its norm on units is the power
\(1+q+\dots+q^{d-1}=(q^d-1)/(q-1)\). A generator therefore has norm
of order \(q-1\), so the norm is onto. The finite-layer
\(H^2(G_E,E^\times)\) vanishes by Proposition~\ref{res:GH11}.
Proposition~\ref{res:GH01} now gives the last equality.
\end{proof}

\begin{example}[A real norm obstruction and two rational-function classes]
\label{res:GH13}
For \(\mathbb C/\mathbb R\), the relative group
\(H^2(\operatorname{Gal}(\mathbb C/\mathbb R),\mathbb C^\times)\)
has order two, with nonzero class represented by a cyclic cocycle
carrying the value \(-1\).
Over \(k=\mathbb F_q(t)\) of characteristic \(p\), \(a=t\) gives a
nonzero Artin--Schreier class. If \(m>1\) divides \(q-1\), the Kummer
class of \(t\) has order \(m\) and gives a cyclic extension of degree
\(m\).
\end{example}
\begin{proof}
The complex norm has image exactly the positive real numbers, by
\(z\bar z>0\) for \(z\ne0\) and positive real square roots. The
first assertion is then Proposition~\ref{res:GH11} with multiplicative
notation. For \(t\), put \(v_\infty(f/g)=\deg g-\deg f\) on
nonzero rational functions in coprime polynomials, and
\(v_\infty(0)=+\infty\). A common denominator shows
\(v_\infty(f+g)\geq\min(v_\infty(f),v_\infty(g))\), with equality
for unequal values because the leading terms cannot cancel.
If \(v_\infty(x)\geq0\), then \(v_\infty(\wp(x))\geq0\); if
\(v_\infty(x)<0\), then \(v_\infty(\wp(x))=p v_\infty(x)\).
Neither gives the value \(-1\) of \(t\), so
Theorem~\ref{res:GH09} applies. Finally \(v_\infty(t^i)=-i\)
and the valuation of any \(m\)-th power is divisible by \(m\).
The class of \(t\) has order \(m\), and \(m\mid q-1\) puts
\(\mu_m\) inside \(k\) by cyclicity of \(\mathbb F_q^\times\)
from Theorem~\ref{res:GH07}. The perfect pairing there gives the
cyclic degree-\(m\) extension.
\end{proof}

\section{Inflation and restriction}

\begin{proposition}[Continuous inflation--restriction in degree one]
\label{res:GH14}
For a profinite group \(G\), a closed normal subgroup \(N\), and a
discrete continuous abelian \(G\)-module \(M\), the sequence
\[
 0\longrightarrow H^1_{\mathrm{cts}}(G/N,M^N)
 \xrightarrow{\mathrm{inf}} H^1_{\mathrm{cts}}(G,M)
 \xrightarrow{\mathrm{res}} H^1_{\mathrm{cts}}(N,M)^{G/N}
\]
is exact at its first two cohomology groups. The action on the last
group comes from \((g\cdot f)(n)=g f(g^{-1}ng)\).
\end{proposition}
\begin{proof}
The displayed action respects the continuous cocycle identity
and sends the coboundary of \(b\in M\) to that of \(gb\).
For any cocycle \(c\) on \(G\), expanding the three factors in
\(g^{-1}ng\) gives
\[
 g c(g^{-1}ng)=c(n)+(n-1)c(g).
\]
The same identity within \(N\), for any cocycle on \(N\), shows
that inner action by elements of \(N\) is trivial on cohomology.
Thus the action factors through \(G/N\), and the identity for
\(c\) on \(G\) places its restriction class among the invariants.

An inflated cocycle vanishes on \(N\). If it is a coboundary
\(g\mapsto(g-1)b\) on \(G\), then \(b\in M^N\), so its original
class was already zero. Conversely a class restricting to zero
has a representative \(c\) vanishing on \(N\), obtained by subtracting
a coboundary whose restriction is the given restriction cocycle.
The cocycle rule gives \(c(gn)=c(g)\) and \(c(ng)=n c(g)\).
Normality gives \(ng\in gN\), so \(n c(g)=c(g)\). Hence \(c\)
descends to a cocycle on \(G/N\) with values in \(M^N\).
It is continuous by the quotient topology. The action of \(G/N\)
on the discrete \(M^N\) is continuous because its point stabilizers
are the images of the open stabilizers in \(G\), each containing
\(N\). Closed normal quotients here are profinite by Result PF02
of \emph{Topics in Group Theory} \cite[PDF pp.~15--16]{mg-groups}.
\end{proof}

\paragraph{Sources.}
The finite-layer and degree-one arithmetic constructions are in
\cite[Sections 2.2, 2.4--2.5]{sharifi} and
\cite[Chapter~I, Sections 2--3 and 5]{serregc}.
For field-theoretic independence, norm and radical expressions,
see \cite[Chapter~5]{milneft}.
\chapter{Embedding problems and regular realizations}
\label{ch:inverse-galois}

An inverse Galois question asks for a Galois extension with a specified
group. An embedding problem also prescribes a quotient extension and
the map to its Galois group. We first treat finite embedding problems,
including both the obstruction to a lift and the extra requirement for
surjectivity. We then pass from rational-function fields to residue
fields by specialization at an explicitly good point.

We use the finite and infinite Galois correspondences of
the earlier field-theory volume~\cite[Results F01, F07, T02,
C04, C06, C08--C10]{mg-aft}. In particular, the absolute Galois
group \(\Gamma_k=\operatorname{Gal}(k^{\mathrm{s}}/k)\) carries
the Krull topology, and finite groups have the discrete topology.

\section{Finite embedding problems}

\begin{definition}[Weak and proper solutions]
\label{res:IG01}
A finite embedding problem for a profinite group \(\Gamma\)
consists of continuous surjections
\[
   \rho\colon\Gamma\twoheadrightarrow Q,\qquad
   \pi\colon B\twoheadrightarrow Q
\]
with \(Q\) and \(B\) finite. Write \(A=\ker\pi\). A weak solution is
a continuous homomorphism \(\lambda\colon\Gamma\to B\) satisfying
\(\pi\lambda=\rho\). It is proper if it is surjective. The problem is
split if \(\pi\) has a homomorphic section and central if \(A\) is
contained in the centre of \(B\).
\end{definition}

\begin{proposition}[The prescribed quotient extension]
\label{res:IG02}
Let \(\Gamma=\Gamma_k\), and let \(\rho\) be restriction to a prescribed
finite Galois extension \(E/k\) in \(k^{\mathrm{s}}\), with the
identification \(Q=\operatorname{Gal}(E/k)\). A weak solution produces
a finite Galois extension \(L/k\) containing \(E\) and an injection
\(\iota\colon\operatorname{Gal}(L/k)\hookrightarrow B\) for which
\(\pi\iota\) is restriction to \(E\). Conversely such an extension
and injection give a weak solution. This is a proper solution exactly
when \(\iota\) is an isomorphism. In particular, when \(Q=1\),
proper solutions are finite Galois realizations of \(B\) over \(k\).
\end{proposition}

\begin{proof}
The kernel of \(\lambda\) is open and normal and is contained in
\(\ker\rho\). The fixed field \(L\) of this kernel is finite and
Galois over \(k\) and contains \(E\). Restriction gives the
isomorphism from the quotient \(\Gamma_k/\ker\lambda\) to
\(\operatorname{Gal}(L/k)\), and \(\lambda\) identifies this
quotient with its image in \(B\). The equation
\(\pi\lambda=\rho\) checks the restriction diagram. Conversely,
compose the surjective restriction to \(L\) with \(\iota\).
\end{proof}

\section{Abelian kernels and properness}

In this section cochains with values in an abelian kernel are written
additively. The coefficients in \(H^i_{\mathrm{cts}}\) have the
discrete topology. This agrees with the cochain convention in
\cite[Results GC01 and GC04--GC05]{mg-groups}.

\begin{lemma}[The class of an abelian extension]
\label{res:IG03}
For an abelian kernel \(A\) in Definition~\ref{res:IG01}, conjugation
by representatives in \(B\) defines an action of \(Q\) on \(A\).
Let \(s\colon Q\to B\) be a set section with \(s(1)=1\).
The factor set \(c(q,r)=s(q)s(r)s(qr)^{-1}\), valued in \(A\),
is a normalized \(2\)-cocycle for this action. Its class
\([c]\in H^2(Q,A)\) is independent of \(s\), and it vanishes
exactly when \(\pi\) splits.
\end{lemma}

\begin{proof}
Representatives of \(q\) differ by an element of \(A\), whose
conjugation on \(A\) is trivial. Also \(s(q)s(r)\) represents \(qr\),
so the resulting conjugations give an action. In the coordinates
\((a,q)\leftrightarrow a\,s(q)\), multiplication has the rule
\[
 (a,q)(b,r)=(a+q\cdot b+c(q,r),\,qr).
\]
Associativity implies
\[
 c(q,r)+c(qr,u)=q\cdot c(r,u)+c(q,ru).
\]
Normalization follows from \(s(1)=1\). Replacing \(s(q)\) by
\(b(q)s(q)\) adds the coboundary
\(\delta b(q,r)=b(q)+q\cdot b(r)-b(qr)\) to \(c\).
If \(c\) is a coboundary, changing the section by the negative of a
cochain representing it makes the new factor set zero. A section
with zero factor set is a homomorphism.
\end{proof}

\begin{theorem}[Obstruction and classification of weak lifts]
\label{res:IG04}
Give \(A\) the \(\Gamma\)-action through \(\rho\).
For an abelian kernel, a weak solution exists exactly when
\[
        \rho^*[c]=0\quad\text{in }H^2_{\mathrm{cts}}(\Gamma,A).
\]
If a weak solution exists, the set of weak solutions modulo
conjugation by \(A\) has a free transitive action of
\(H^1_{\mathrm{cts}}(\Gamma,A)\).
For a central kernel the action on \(A\) is trivial.
\end{theorem}

\begin{proof}
A continuous map lifting \(\rho\) has the unique form
\(\lambda(\gamma)=z(\gamma)s(\rho\gamma)\) with a continuous
map \(z\colon\Gamma\to A\). By Lemma~\ref{res:IG03}, it is a
homomorphism exactly when
\[
   z(\gamma)+\gamma\cdot z(\eta)-z(\gamma\eta)
           =-c(\rho\gamma,\rho\eta).
\]
The equation at \((1,1)\) forces \(z(1)=0\).
This gives the obstruction. Once one such \(z_0\) is fixed, all
solutions are \(z_0+w\) for continuous \(1\)-cocycles \(w\).
Conjugation by \(a\in A\) adds the cochain
\(\gamma\mapsto a-\gamma\cdot a\) to \(z\); these cochains are
exactly the \(1\)-coboundaries. This proves the claimed action.
\end{proof}

\begin{corollary}[Detection on finite stages]
\label{res:IG05}
Let \(U\) range over the open normal subgroups of \(\Gamma\) contained
in \(\ker\rho\), and write \(\rho_U\colon\Gamma/U\to Q\) for
the induced map. For abelian kernel \(A\), a weak solution exists
exactly when \(\rho_U^*[c]\) vanishes at one finite stage
\(H^2(\Gamma/U,A)\). Inflation along
\(\Gamma/V\to\Gamma/U\), for \(V\leq U\), carries the class
for \(U\) to the class for \(V\). In an absolute Galois group these
are the stages obtained by restriction to the finite Galois
extensions containing \(E\).
\end{corollary}

\begin{proof}
The proof of Theorem~\ref{res:IG04} applied to the finite group
\(\Gamma/U\) gives the equivalence between vanishing at this
stage and a weak lift from this stage. Inflation acts by
precomposition on each cochain, giving the asserted transition.
Every continuous lift to the finite group \(B\) factors through
the finite quotient by its open normal kernel, which is contained
in \(\ker\rho\). For \(\Gamma_k\), open normal kernels are exactly
the kernels of restriction to finite Galois subextensions, by
the earlier correspondence~\cite[Results C06 and C08--C09]{mg-aft}.
\end{proof}

\begin{proposition}[Proper solutions of split problems]
\label{res:IG06}
Put \(N=\ker\rho\). A weak solution of any finite embedding problem
is proper if and only if \(\lambda(N)=A\).
Suppose \(\pi\) splits and identify \(B=A\rtimes Q\) using a
homomorphic section. Weak solutions are the maps
\[
 \lambda(\gamma)=(z(\gamma),\rho\gamma),\qquad
 z(\gamma\eta)=z(\gamma)\bigl(\rho\gamma\cdot z(\eta)\bigr),
\]
for continuous \(z\colon\Gamma\to A\); \(A\) may be nonabelian
in this assertion. The restriction \(z|_N\) is a homomorphism,
and this weak solution is proper exactly when \(z(N)=A\).
\end{proposition}

\begin{proof}
The image \(H\) of any weak solution maps onto \(Q\).
The equality \(H\cap A=\lambda(N)\) follows by applying \(\pi\)
and using \(\pi\lambda=\rho\). If \(H\cap A=A\), then for
any \(b\in B\) choose \(h\in H\) over \(\pi(b)\) and write
\(b=(bh^{-1})h\in H\). The converse is immediate.
For a split problem the displayed equation for \(z\) is exactly
the multiplication condition in the semidirect product. On \(N\),
the action through \(\rho\) is trivial and \(z\) is the
restriction of \(\lambda\) with values in \(A\).
\end{proof}

\begin{example}[A central obstruction and a split weak solution]
\label{res:IG07}
Take \(\Gamma=Q=C_2\) and \(\rho\) the identity. The epimorphism
\(C_4\to C_2\) has no weak solution. Indeed the two possible
images of the generator both have order four. The lift from the
set section has \(c(\bar 1,\bar 1)=2\in\{0,2\}\subset C_4\).
For a normalized cochain \(z\) with values in this trivial-action
kernel, \(\delta z(\bar 1,\bar 1)=2z(\bar 1)=0\). Thus the
factor set represents a nonzero obstruction in the notation of
Theorem~\ref{res:IG04}. In comparison, projection
\(C_2\times C_2\to C_2\) admits the section in the first
factor. Its embedding problem has weak solutions, and
Proposition~\ref{res:IG06} rules out a proper one since \(N=1\).
\end{example}

\begin{proposition}[Finite-field embedding problems]
\label{res:IG08}
Every finite embedding problem for the absolute Galois group of
a finite field has a weak solution. It has a proper solution
exactly when \(B\) is cyclic.
\end{proposition}

\begin{proof}
The absolute Galois group is \(\widehat{\mathbb Z}\) with
arithmetic Frobenius the topological generator
\cite[Results E04, E06 and E07]{mg-aft}. The element
\(q=\rho(1)\) therefore generates \(Q\). Any \(b\in\pi^{-1}(q)\)
of order \(m\) defines a weak lift by reducing
\(\widehat{\mathbb Z}\) modulo \(m\) and sending \(1\) to \(b\).
The composite with \(\pi\) agrees with \(\rho\) on the dense
integer subgroup and hence everywhere. Since the image of
\(\widehat{\mathbb Z}\) in any finite group is generated by
the image of \(1\), properness requires \(B\) cyclic.

Conversely, suppose \(B=\langle x\rangle\) has order \(m\) and
\(Q\) has order \(n\mid m\). We have
\(q=\pi(x)^a\) for some \(a\) prime to \(n\).
Choose \(a'\equiv a\pmod n\) and \(a'\equiv1\pmod p\) for
each prime \(p\mid m\) with \(p\nmid n\), using the Chinese
remainder construction. Then \(a'\) is prime to \(m\), and
\(x^{a'}\) is a generator of \(B\) lifting \(q\).
\end{proof}

\section{Specialization at a good point}

Let \(k\) be any field and \(u\) an indeterminate. For a field
\(\ell\) containing \(k\), write \(v_u\) for the order of a nonzero
Laurent series over \(\ell\). The leading terms of nonzero
series multiply, so orders add under multiplication, and the order of
a sum is at least the lesser order. The residue of a series of
nonnegative order is its constant term.

\begin{lemma}[Simple-root lifting in power series]
\label{res:IG09}
If \(f\in k[[u]][X]\), \(\ell/k\) is a field extension, and
\(r\in\ell\) is a simple root of the reduction \(\overline f\),
there is exactly one \(\alpha\in\ell[[u]]\) with
\(\alpha(0)=r\) and \(f(\alpha)=0\).
\end{lemma}

\begin{proof}
Begin with \(\alpha_1=r\), which solves the equation modulo \(u\).
Given a polynomial \(\alpha_j\equiv r\pmod u\) with
\(f(\alpha_j)\equiv0\pmod{u^j}\) for \(j\geq1\),
every extension of its coefficients through degree \(j\) has
the form \(\alpha_j+b u^j\) modulo \(u^{j+1}\).
Expansion gives
\[
 f(\alpha_j+b u^j)\equiv f(\alpha_j)+b u^j\overline f'(r)
          \pmod{u^{j+1}},
\]
since \(2j\geq j+1\). The coefficient \(b\in\ell\) that kills the
degree-\(j\) term is unique because \(\overline f'(r)\ne0\).
The successive coefficients define a series \(\alpha\).
Its evaluation in \(f\) is zero modulo every \(u^j\), so it
is zero. The same coefficient recursion proves uniqueness.
\end{proof}

\begin{theorem}[Good specialization as a decomposition subgroup]
\label{res:IG10}
Let \(f(t,X)\in k(t)[X]\) be monic, with coefficients regular at
\(t=a\in k\), and assume \(f(a,X)\) is separable. Let \(L\) be a
splitting field over \(k(t)\) and \(\ell\) a splitting field of
\(f(a,X)\) over \(k\). There is an embedding
\(L\hookrightarrow\ell((u))\) extending \(t\mapsto a+u\)
which matches the generic roots with their distinct reductions.
For the valuation \(w\) induced by \(v_u\), its residue field is
\(\ell\), its value group is \(\mathbb Z\) with \(w(t-a)=1\),
and the map
\[
 \{\sigma\in\operatorname{Gal}(L/k(t)):
           w(\sigma x)=w(x)\ \text{for every }x\in L^\times\}
       \longrightarrow \operatorname{Gal}(\ell/k)
\]
given by reduction is an isomorphism. In particular the subgroup
stabilizing \(w\) has trivial inertia kernel, and the specialized
Galois group embeds into the generic Galois group up to conjugacy.
\end{theorem}

\begin{proof}
The substitution \(t=a+u\) gives the indicated embedding of
\(k(t)\) into \(k((u))\). A coefficient regular at \(a\) can be
written with a denominator nonzero at \(a\), whose image in
\(k[[u]]\) has an inverse by coefficient recursion.
By Lemma~\ref{res:IG09}, each root of \(f(a,X)\) lifts to a root
of \(f(a+u,X)\) in \(\ell[[u]]\). The lifted roots are distinct
by their reductions. Since the polynomial is monic of unchanged
degree, the lifts are all of its roots. They generate a splitting
field over the embedded \(k(t)\); uniqueness of the splitting field
gives the required embedding. Write \(L\) for its image.

The residue field of \(L\) lies in \(\ell\) and contains every
specialized root, so equals \(\ell\). The value group lies in
\(\mathbb Z\) and contains \(1=w(u)\). An automorphism preserving
\(w\) acts on the residue field. If it induces the identity, it
fixes the reduction of each lifted root. The reductions are distinct,
and the automorphism permutes these roots, so it fixes all of them
and is the identity. For surjectivity, extend any
\(\tau\in\operatorname{Gal}(\ell/k)\) coefficientwise to
\(\ell((u))\), fixing \(u\). This automorphism preserves order.
It takes the lift of \(r\) to the lift of \(\tau(r)\) by uniqueness,
and hence preserves \(L\). Its restriction induces \(\tau\) on
the residue field. Finally, a different embedding into
\(\ell((u))\) over the same \(t\mapsto a+u\) has the same image,
since its roots must be the same full set of lifted roots.
Thus it changes the induced valuation by an automorphism of \(L\),
which conjugates the displayed subgroup.
\end{proof}

\begin{corollary}[Specialization in a tower]
\label{res:IG11}
In Theorem~\ref{res:IG10}, let \(L_0\subseteq L\) be a splitting
field over \(k(t)\) of another monic polynomial \(f_0(t,X)\)
regular at \(a\) with separable reduction. Under the valuation
induced from \(L\), the residue field \(\ell_0\) of \(L_0\) is
the splitting field of \(f_0(a,X)\) inside \(\ell\).
Under the decomposition group isomorphisms of
Theorem~\ref{res:IG10}, restriction from the subgroup for \(L\)
to the subgroup for \(L_0\) is exactly the surjective map
\(\operatorname{Gal}(\ell/k)\to\operatorname{Gal}(\ell_0/k)\).
\end{corollary}

\begin{proof}
Every root of \(f_0(a+u,X)\) in \(L_0\) is integral for the
series valuation. If a root had negative order, its leading monic
term in the polynomial equation would have strictly lower order
than each of the other terms, which cannot sum to zero.
Reducing the monic factorization gives all the roots of \(f_0(a,X)\)
in \(\ell\), and they are distinct by separability. Let \(\ell_0\)
be the field they generate. The same roots lift in
\(\ell_0[[u]]\) by Lemma~\ref{res:IG09}, and uniqueness of the
lifts inside \(\ell[[u]]\) identifies these with the original
roots. Consequently \(L_0\) embeds into \(\ell_0((u))\), and its
residue field is exactly \(\ell_0\). Reduction commutes with
restriction to \(L_0\). The residue restriction is surjective
by finite Galois restriction, and the two vertical identifications
are those of Theorem~\ref{res:IG10}.
\end{proof}

\section{Regular realizations and arithmetic examples}

A finite Galois realization \(L/k(t)\) is called regular over
\(k\) when \(k\) is algebraically closed in \(L\): every element
of \(L\) algebraic over \(k\) lies in \(k\).
Here the finite extension of \(k(t)\) is separable already.
The criterion below works over imperfect fields as well.

\begin{proposition}[Regularity and the geometric group]
\label{res:IG12}
For a finite Galois extension \(L/k(t)\), work in a common algebraic
closure of \(k(t)\) and write \(\overline k\) for its algebraic
closure of constants. Restriction embeds
\(\operatorname{Gal}(L\overline k(t)/\overline k(t))\) into
\(\operatorname{Gal}(L/k(t))\). It is onto exactly when \(L/k(t)\)
is regular over \(k\).
\end{proposition}

\begin{proof}
Let \(M=L\cap\overline k(t)\). Restriction identifies the geometric
group with \(\operatorname{Gal}(L/M)\) by Galois base change
\cite[Result C10]{mg-aft}. We show \(M=C(t)\) for a finite
extension \(C/k\) inside \(L\cap\overline k\).

In characteristic \(p>0\), each member of \(\overline k(t)\)
has some \(p\)-power in \(k^{\mathrm{s}}(t)\): take an exponent
making all its finitely many algebraic coefficients separable over
\(k\). For an algebraic constant \(\beta\), write its monic minimal
polynomial as \(g(X^{p^r})\) with \(g'\ne0\) by collecting successive
\(p\)-th powers in the exponents. Irreducibility of \(g\) follows
from that of the original polynomial, so \(\beta^{p^r}\) is
separable. Any \(x\in M\) is separable over \(k(t)\), and hence
separable over \(k^{\mathrm{s}}(t)\), since its minimal polynomial
over the larger field divides a separable polynomial. It is also
purely inseparable over \(k^{\mathrm{s}}(t)\), so it belongs there.
In characteristic zero \(k^{\mathrm{s}}=\overline k\) already.
Since \(M/k(t)\) is finite, its finitely many generators have
their coefficients in some finite Galois extension \(k'/k\)
inside \(k^{\mathrm{s}}\), using Result F03 of
\cite{mg-aft}.

The extension \(k'(t)/k(t)\) has group \(\operatorname{Gal}(k'/k)\)
acting on coefficients. Indeed any constant basis over \(k\)
remains independent over \(k(t)\) by clearing denominators and
comparing powers of \(t\). The resulting span over \(k(t)\)
is a finite-dimensional domain containing \(k'\) and \(k(t)\).
Multiplication by a nonzero element is injective and hence
surjective, so this span is a field and equals \(k'(t)\).
Thus the coefficient actions already supply as many automorphisms
as the extension degree. Let \(H\) be the subgroup fixing \(M\),
and put \(C=(k')^H\). Then \(C(t)\subseteq M\); the same
basis argument gives \([k'(t):C(t)]=[k':C]=|H|\), which is
also \([k'(t):M]\) by finite Galois correspondence. Hence
\(M=C(t)\), with \(C\subseteq L\cap\overline k\).
If \(L\cap\overline k=k\), this forces \(M=k(t)\).
Conversely if \(M=k(t)\), an algebraic constant in \(L\) is
a constant rational function in \(k(t)\), thus lies in \(k\)
by comparison of coefficients after clearing the denominator.
\end{proof}

\begin{theorem}[A regular \(S_3\)-family over \(\mathbb Q\)]
\label{res:IG13}
The splitting field of \(X^3-X-t\) over \(\mathbb Q(t)\)
is regular over \(\mathbb Q\) with group \(S_3\).
Every odd integer \(a\) gives a good specialization with
Galois group \(S_3\) over \(\mathbb Q\).
\end{theorem}

\begin{proof}
There is no root in \(\overline{\mathbb Q}(t)\).
For a putative root \(P(t)/Q(t)\) in lowest terms, the equation
\(P^3-PQ^2=tQ^3\) forces \(Q\mid P^3\) in
\(\overline{\mathbb Q}[t]\). Thus \(Q\) is constant, and
a nonconstant polynomial \(h\) with \(h^3-h=t\) is impossible
by degree. A reducible cubic over a field would have a root.

For a monic cubic with roots \(r_1,r_2,r_3\) counted with multiplicity,
a square root of its discriminant is
\(d=\prod_{i<j}(r_i-r_j)\). Permuting the roots acts on \(d\)
by the sign. For \(f=X^3-X-t\) one has
\[
 d^2=-\prod_i f'(r_i)
     =-27f(b)f(-b)=4-27t^2,\qquad b^2=1/3 .
\]
The first equality follows by taking each ordered pair of
distinct indices in the derivative product; the second follows
from \(f'(X)=3(X-b)(X+b)\). The displayed polynomial is
nonzero, so the cubic is separable, and it has two distinct
simple zeros in \(\overline{\mathbb Q}\), so it is no square
in \(\overline{\mathbb Q}(t)\). For an irreducible separable
cubic in characteristic zero, its Galois group is a transitive
subgroup of \(S_3\). It has order divisible by three and it
lies in \(A_3\) exactly when \(d\) lies in the ground field.
Indeed this is the fixed-field criterion for \(d\), and
\(d^2\) is a square in the ground field exactly when \(d\) is
in it. Thus a nonsquare discriminant gives order six.
This proves that the geometric group is \(S_3\); the arithmetic
group also acts faithfully on the three roots and contains the
geometric group. Regularity follows by Proposition~\ref{res:IG12}.

If \(a\) is odd, the specialization \(X^3-X-a\) reduces to
\(X^3+X+1\) modulo \(2\), which has neither \(0\) nor \(1\)
as a root. A rational root of a monic integer cubic is an
integer, as its denominator in lowest terms divides the cube
of its numerator. There is no such root here. Its discriminant
\(4-27a^2\) is negative, so it is nonzero and nonsquare over
\(\mathbb Q\). The same transitivity and sign argument gives
group \(S_3\). Theorem~\ref{res:IG10} identifies this specialized
group with the full generic group.
\end{proof}

\begin{example}[A regular cyclic family]
\label{res:IG14}
Let \(n\geq2\) be invertible in \(k\), and suppose \(k\) contains a
primitive \(n\)-th root of unity \(\zeta\). The
splitting field of \(X^n-t\) over \(k(t)\) is a regular realization
of \(C_n\). The exponent of the factor \(t\) in a rational
function, with poles counted negatively, is additive under
multiplication. Hence \(t^i\) can be an \(n\)-th power in \(k(t)\)
only if \(n\mid i\), and the class of \(t\) has order \(n\).
Kummer duality (Theorem~\ref{res:GH07}) gives the cyclic group.
The same exponent argument over \(\overline k(t)\) gives the
full geometric group, so regularity follows from
Proposition~\ref{res:IG12}. For instance \(n=2\) over
\(\mathbb Q\) gives full group at each positive prime \(t=p\)
by Theorem~\ref{res:IG10}: the specialization has distinct roots
and \(p\) is nonsquare by prime factorization.
\end{example}

\begin{lemma}[Fresh primes for prescribed cyclic factors]
\label{res:IG15}
Given \(n\geq2\) and a finite set of primes, there is a prime
\(p\) outside the set with \(n\mid p-1\).
\end{lemma}

\begin{proof}
Use the monic integral cyclotomic polynomials and their product
formula from \cite[Result K07.15.cyclotomic\_polynomials]{mg-d07}.
In that formula \(\Phi_1(0)=-1\); evaluation at zero and induction
on \(n\) give \(\Phi_n(0)=1\) for \(n>1\).
Choose an integer \(b\geq3\) divisible by \(n\) and all the
excluded primes. Since each primitive complex \(n\)-th root
\(\xi\) has modulus \(1\), each factor \(|b-\xi|\) is at least
\(b-1\). Thus \(|\Phi_n(b)|>1\). Choose a prime divisor \(p\)
of \(\Phi_n(b)\); it does not divide \(b\), since
\(\Phi_n(b)\equiv 1\) modulo every prime dividing \(b\).
In particular \(p\nmid n\). Modulo \(p\) the polynomial
\(X^n-1\) has no repeated root, since its derivative is
\(nX^{n-1}\). The factors \(\Phi_d\) for \(d\mid n\) therefore
have disjoint root sets. The class of \(b\) has \(n\)-th power
one; if its order were \(d<n\) with \(d\mid n\), it would
also be a root of a factor of \(X^d-1\), a contradiction.
Hence it has order \(n\) in \(\mathbb F_p^\times\), and
\(n\mid p-1\).
\end{proof}

\begin{theorem}[Abelian realizations disjoint from a prescribed field]
\label{res:IG16}
Every finite abelian group \(A\) occurs as a Galois group over
\(\mathbb Q\). More precisely, for any finite Galois extension
\(E/\mathbb Q\) there is a finite Galois extension \(F/\mathbb Q\)
with group \(A\) and \(F\cap E=\mathbb Q\). Consequently the
split embedding problem
\[
   A\times\operatorname{Gal}(E/\mathbb Q)
             \longrightarrow \operatorname{Gal}(E/\mathbb Q)
\]
over the restriction to \(E\) has a proper solution through \(EF\).
\end{theorem}

\begin{proof}
Write \(A=\prod_{i=1}^r C_{n_i}\) with \(n_i\geq2\),
using \cite[Result K07.07.invariant\_factor\_existence]{mg-d07}.
The case \(r=0\) takes \(F=\mathbb Q\). Choose any finite
number \(s\) of disjoint prime lists \(p_{j,i}\), with
\(n_i\mid p_{j,i}-1\), by Lemma~\ref{res:IG15}.
Let \(m\) be the product of all selected primes. The cyclotomic
group of \(\mathbb Q(\zeta_m)\) is
\((\mathbb Z/m\mathbb Z)^\times\) by
\cite[Result K07.15.cyclotomic\_galois]{mg-d07}.
By the Chinese remainder isomorphism it is the product of
the groups \(\mathbb F_{p_{j,i}}^\times\), each cyclic of
order \(p_{j,i}-1\) by
\cite[Result K07.13.multiplicative\_cyclic]{mg-d07}.
For list \(j\), project to its factors and take their respective
quotients of order \(n_i\), giving a surjection onto a copy
of \(A\). The fixed field \(F_j\) of its kernel is Galois
with group \(A\). For two distinct lists the joint map onto
\(A\times A\) is surjective; thus
\([F_jF_h:\mathbb Q]=|A|^2\) for \(j\ne h\), by the
fixed-field correspondence. Galois base change
\cite[Result C10]{mg-aft} now gives \(F_j\cap F_h=\mathbb Q\).

The finite Galois extension \(E/\mathbb Q\) has only finitely
many intermediate fields, since its Galois group has finitely
many subgroups. Take \(s\) greater than this number. If all
intersections \(F_j\cap E\) were nontrivial over \(\mathbb Q\),
they would be pairwise different because the \(F_j\) have
trivial pairwise intersection. This is impossible, so some
\(F=F_j\) has \(F\cap E=\mathbb Q\). Restriction from
\(\operatorname{Gal}(EF/\mathbb Q)\) to the product for \(F\)
and \(E\) is injective since the two fields generate \(EF\);
it is surjective by the degree formula from Result C10.
The resulting isomorphism has second projection equal to
restriction to \(E\). Proposition~\ref{res:IG02} gives the
proper solution.
\end{proof}
\backmatter
\begingroup
\small
\begin{thebibliography}{99}
\phantomsection
\addcontentsline{toc}{chapter}{Bibliography}
\renewcommand{\runningchapter}{Bibliography}
\setlength{\itemsep}{4pt}

\bibitem{debes}
P. D\`ebes.
\emph{Arithm\'etique des rev\^etements de la droite}.
Course notes, version 3. Chapter 5 treats the inverse Galois problem
and specialization.

\bibitem{haranjarden}
D. Haran and M. Jarden.
\emph{Regular Split Embedding Problems over Complete Valued Fields}.
Forum Mathematicum 10, 1998.

\bibitem{mallematzat}
G. Malle and B.~H. Matzat.
\emph{Inverse Galois Theory}.
Second edition, Springer, 2018.

\bibitem{milneft}
J.~S. Milne.
\emph{Fields and Galois Theory}.
Author's course notes, version 5.10.

\bibitem{roques}
J. Roques.
\emph{An Introduction to Difference Galois Theory}.
CRM summer school, Montr\'eal, 2016.

\bibitem{serregc}
J.-P. Serre.
\emph{Galois Cohomology}.
Translated by P. Ion, Springer, English first printing, 1997.

\bibitem{sharifi}
R. Sharifi.
\emph{Group and Galois Cohomology}.
UCLA course notes. The archived version of Chapter 2 contains the
treatment through Section 2.5.

\bibitem{singer}
M.~F. Singer.
\emph{Introduction to the Galois Theory of Linear Differential Equations}.
Expanded notes for the 2006 LMS Short Instructional Course on
Differential Algebra and Related Topics, arXiv:0712.4124v2, 2008.

\bibitem{vps-differential}
M. van der Put and M.~F. Singer.
\emph{Galois Theory of Linear Differential Equations}.
Author's revised manuscript, July 2002. The book appeared in
Grundlehren der mathematischen Wissenschaften 328, Springer, 2003.

\bibitem{vps-difference}
M. van der Put and M.~F. Singer.
\emph{Galois Theory of Difference Equations}.
Lecture Notes in Mathematics 1666, Springer, 1997.

\bibitem{mg-aft}
MathGenome.
\emph{Infinite Galois Theory}.
ALG-04 core notes, 22-page revision, 2026.
Results F01--F07 concern finite subextensions, T01--T08 the Krull
topology, C01--C12 the Galois correspondence, and E01--E09 finite fields.

\bibitem{mg-groups}
MathGenome.
\emph{Topics in Group Theory}.
ALG-02 deepening notes, 31-page revision, 2026.
Chapter 3 covers profinite groups. Chapter 4 covers continuous
cohomology, extensions and tree actions.

\bibitem{mg-d07}
MathGenome.
\emph{Abstract Algebra}.
D07, revision archived in MathGenome v0.25.0.

\bibitem{mg-algebra-core}
MathGenome.
\emph{Algebra Core}.
P1-03, revision archived in MathGenome v0.25.0.
Theorem 10.3, \texttt{p103:ALG-10-C01-R03}, is the long exact sequence
of a short exact sequence of complexes.

\bibitem{mg-ha}
MathGenome.
\emph{Advanced Homological Algebra}.
ALG-10, 62-page revision archived in MathGenome v0.25.0.

\bibitem{mg-rm03}
MathGenome.
\emph{Rings, Modules, and Associative Algebras}.
ALG-03 core notes, 19-page revision, 2026.

\bibitem{mg-rm03d}
MathGenome.
\emph{Topics in Rings, Modules, and Associative Algebras}.
ALG-03 deepening notes, 38-page revision, 2026.
Chapter 4 studies differential and difference operators.

\end{thebibliography}
\endgroup
\end{document}
